% Options for packages loaded elsewhere
% Options for packages loaded elsewhere
\PassOptionsToPackage{unicode}{hyperref}
\PassOptionsToPackage{hyphens}{url}
\PassOptionsToPackage{dvipsnames,svgnames,x11names}{xcolor}
%
\documentclass[
  turkish,
  a4paper,
  DIV=11,
  numbers=noendperiod]{scrreprt}
\usepackage{xcolor}
\usepackage[margin=2.5cm]{geometry}
\usepackage{amsmath,amssymb}
\setcounter{secnumdepth}{5}
\usepackage{iftex}
\ifPDFTeX
  \usepackage[T1]{fontenc}
  \usepackage[utf8]{inputenc}
  \usepackage{textcomp} % provide euro and other symbols
\else % if luatex or xetex
  \usepackage{unicode-math} % this also loads fontspec
  \defaultfontfeatures{Scale=MatchLowercase}
  \defaultfontfeatures[\rmfamily]{Ligatures=TeX,Scale=1}
\fi
\usepackage{lmodern}
\ifPDFTeX\else
  % xetex/luatex font selection
\fi
% Use upquote if available, for straight quotes in verbatim environments
\IfFileExists{upquote.sty}{\usepackage{upquote}}{}
\IfFileExists{microtype.sty}{% use microtype if available
  \usepackage[]{microtype}
  \UseMicrotypeSet[protrusion]{basicmath} % disable protrusion for tt fonts
}{}
\makeatletter
\@ifundefined{KOMAClassName}{% if non-KOMA class
  \IfFileExists{parskip.sty}{%
    \usepackage{parskip}
  }{% else
    \setlength{\parindent}{0pt}
    \setlength{\parskip}{6pt plus 2pt minus 1pt}}
}{% if KOMA class
  \KOMAoptions{parskip=half}}
\makeatother
% Make \paragraph and \subparagraph free-standing
\makeatletter
\ifx\paragraph\undefined\else
  \let\oldparagraph\paragraph
  \renewcommand{\paragraph}{
    \@ifstar
      \xxxParagraphStar
      \xxxParagraphNoStar
  }
  \newcommand{\xxxParagraphStar}[1]{\oldparagraph*{#1}\mbox{}}
  \newcommand{\xxxParagraphNoStar}[1]{\oldparagraph{#1}\mbox{}}
\fi
\ifx\subparagraph\undefined\else
  \let\oldsubparagraph\subparagraph
  \renewcommand{\subparagraph}{
    \@ifstar
      \xxxSubParagraphStar
      \xxxSubParagraphNoStar
  }
  \newcommand{\xxxSubParagraphStar}[1]{\oldsubparagraph*{#1}\mbox{}}
  \newcommand{\xxxSubParagraphNoStar}[1]{\oldsubparagraph{#1}\mbox{}}
\fi
\makeatother


\usepackage{longtable,booktabs,array}
\newcounter{none} % for unnumbered tables
\usepackage{calc} % for calculating minipage widths
% Correct order of tables after \paragraph or \subparagraph
\usepackage{etoolbox}
\makeatletter
\patchcmd\longtable{\par}{\if@noskipsec\mbox{}\fi\par}{}{}
\makeatother
% Allow footnotes in longtable head/foot
\IfFileExists{footnotehyper.sty}{\usepackage{footnotehyper}}{\usepackage{footnote}}
\makesavenoteenv{longtable}
\usepackage{graphicx}
\makeatletter
\newsavebox\pandoc@box
\newcommand*\pandocbounded[1]{% scales image to fit in text height/width
  \sbox\pandoc@box{#1}%
  \Gscale@div\@tempa{\textheight}{\dimexpr\ht\pandoc@box+\dp\pandoc@box\relax}%
  \Gscale@div\@tempb{\linewidth}{\wd\pandoc@box}%
  \ifdim\@tempb\p@<\@tempa\p@\let\@tempa\@tempb\fi% select the smaller of both
  \ifdim\@tempa\p@<\p@\scalebox{\@tempa}{\usebox\pandoc@box}%
  \else\usebox{\pandoc@box}%
  \fi%
}
% Set default figure placement to htbp
\def\fps@figure{htbp}
\makeatother


% definitions for citeproc citations
\NewDocumentCommand\citeproctext{}{}
\NewDocumentCommand\citeproc{mm}{%
  \begingroup\def\citeproctext{#2}\cite{#1}\endgroup}
\makeatletter
 % allow citations to break across lines
 \let\@cite@ofmt\@firstofone
 % avoid brackets around text for \cite:
 \def\@biblabel#1{}
 \def\@cite#1#2{{#1\if@tempswa , #2\fi}}
\makeatother
\newlength{\cslhangindent}
\setlength{\cslhangindent}{1.5em}
\newlength{\csllabelwidth}
\setlength{\csllabelwidth}{3em}
\newenvironment{CSLReferences}[2] % #1 hanging-indent, #2 entry-spacing
 {\begin{list}{}{%
  \setlength{\itemindent}{0pt}
  \setlength{\leftmargin}{0pt}
  \setlength{\parsep}{0pt}
  % turn on hanging indent if param 1 is 1
  \ifodd #1
   \setlength{\leftmargin}{\cslhangindent}
   \setlength{\itemindent}{-1\cslhangindent}
  \fi
  % set entry spacing
  \setlength{\itemsep}{#2\baselineskip}}}
 {\end{list}}
\usepackage{calc}
\newcommand{\CSLBlock}[1]{\hfill\break\parbox[t]{\linewidth}{\strut\ignorespaces#1\strut}}
\newcommand{\CSLLeftMargin}[1]{\parbox[t]{\csllabelwidth}{\strut#1\strut}}
\newcommand{\CSLRightInline}[1]{\parbox[t]{\linewidth - \csllabelwidth}{\strut#1\strut}}
\newcommand{\CSLIndent}[1]{\hspace{\cslhangindent}#1}

\ifLuaTeX
\usepackage[bidi=basic,shorthands=off]{babel}
\else
\usepackage[bidi=default,shorthands=off]{babel}
\fi
\ifLuaTeX
  \usepackage{selnolig} % disable illegal ligatures
\fi


\setlength{\emergencystretch}{3em} % prevent overfull lines

\providecommand{\tightlist}{%
  \setlength{\itemsep}{0pt}\setlength{\parskip}{0pt}}



 


\KOMAoption{captions}{tableheading}
\makeatletter
\@ifpackageloaded{tcolorbox}{}{\usepackage[skins,breakable]{tcolorbox}}
\@ifpackageloaded{fontawesome5}{}{\usepackage{fontawesome5}}
\definecolor{quarto-callout-color}{HTML}{909090}
\definecolor{quarto-callout-note-color}{HTML}{0758E5}
\definecolor{quarto-callout-important-color}{HTML}{CC1914}
\definecolor{quarto-callout-warning-color}{HTML}{EB9113}
\definecolor{quarto-callout-tip-color}{HTML}{00A047}
\definecolor{quarto-callout-caution-color}{HTML}{FC5300}
\definecolor{quarto-callout-color-frame}{HTML}{acacac}
\definecolor{quarto-callout-note-color-frame}{HTML}{4582ec}
\definecolor{quarto-callout-important-color-frame}{HTML}{d9534f}
\definecolor{quarto-callout-warning-color-frame}{HTML}{f0ad4e}
\definecolor{quarto-callout-tip-color-frame}{HTML}{02b875}
\definecolor{quarto-callout-caution-color-frame}{HTML}{fd7e14}
\makeatother
\makeatletter
\@ifpackageloaded{caption}{}{\usepackage{caption}}
\AtBeginDocument{%
\ifdefined\contentsname
  \renewcommand*\contentsname{İçindekiler}
\else
  \newcommand\contentsname{İçindekiler}
\fi
\ifdefined\listfigurename
  \renewcommand*\listfigurename{ŞEKİLLER DİZİNİ}
\else
  \newcommand\listfigurename{ŞEKİLLER DİZİNİ}
\fi
\ifdefined\listtablename
  \renewcommand*\listtablename{TABLOLAR DİZİNİ}
\else
  \newcommand\listtablename{TABLOLAR DİZİNİ}
\fi
\ifdefined\figurename
  \renewcommand*\figurename{Şekil}
\else
  \newcommand\figurename{Şekil}
\fi
\ifdefined\tablename
  \renewcommand*\tablename{Tablo}
\else
  \newcommand\tablename{Tablo}
\fi
}
\@ifpackageloaded{float}{}{\usepackage{float}}
\floatstyle{ruled}
\@ifundefined{c@chapter}{\newfloat{codelisting}{h}{lop}}{\newfloat{codelisting}{h}{lop}[chapter]}
\floatname{codelisting}{Listeleme}
\newcommand*\listoflistings{\listof{codelisting}{İlan Listesi}}
\makeatother
\makeatletter
\makeatother
\makeatletter
\@ifpackageloaded{caption}{}{\usepackage{caption}}
\@ifpackageloaded{subcaption}{}{\usepackage{subcaption}}
\makeatother
\usepackage{bookmark}
\IfFileExists{xurl.sty}{\usepackage{xurl}}{} % add URL line breaks if available
\urlstyle{same}
\hypersetup{
  pdftitle={T1DM Niteliksel Kol --- Kanonik Bilimsel Sonuç Belgesi},
  pdfauthor={Niteliksel araştırma kolu (OM birinci araştırmacı · BA critical friend)},
  pdflang={tr},
  colorlinks=true,
  linkcolor={blue},
  filecolor={Maroon},
  citecolor={Blue},
  urlcolor={Blue},
  pdfcreator={LaTeX via pandoc}}


\title{T1DM Niteliksel Kol --- Kanonik Bilimsel Sonuç Belgesi}
\usepackage{etoolbox}
\makeatletter
\providecommand{\subtitle}[1]{% add subtitle to \maketitle
  \apptocmd{\@title}{\par {\large #1 \par}}{}{}
}
\makeatother
\subtitle{Tip 1 Diyabet \& Ebeveynlik Tutumu doktora tezinin niteliksel
kolu: nihai analiz ve sonuçlar}
\author{Niteliksel araştırma kolu (OM birinci araştırmacı · BA critical
friend)}
\date{2026-07-11}
\begin{document}
\maketitle

\renewcommand*\contentsname{İçindekiler}
{
\setcounter{tocdepth}{2}
\tableofcontents
}

\begin{tcolorbox}[enhanced jigsaw, arc=.35mm, bottomrule=.15mm, breakable, colback=white, colframe=quarto-callout-note-color-frame, left=2mm, leftrule=.75mm, opacityback=0, rightrule=.15mm, toprule=.15mm]
\begin{minipage}[t]{5.5mm}
\textcolor{quarto-callout-note-color}{\faInfo}
\end{minipage}%
\begin{minipage}[t]{\textwidth - 5.5mm}

\textbf{Belge statüsü.} Bu Quarto belgesi niteliksel kolun \textbf{tek
kanonik bilimsel sonuç kaynağıdır}. Repo kökündeki
\texttt{niteliksel/qualitative\_canonical\_results\_report.md}, bu
dosyanın mekanik Markdown kopyası olarak tutulur; ikinci bir içerik
otoritesi değildir. Kanonik sayılar, tema mimarisi, bulgular,
metodoloji, güvenilirlik/değerlendirme ve raporlama kesitleri bu dosyada
yönetilir.

\end{minipage}%
\end{tcolorbox}

\begin{tcolorbox}[enhanced jigsaw, arc=.35mm, bottomrule=.15mm, breakable, colback=white, colframe=quarto-callout-important-color-frame, left=2mm, leftrule=.75mm, opacityback=0, rightrule=.15mm, toprule=.15mm]
\begin{minipage}[t]{5.5mm}
\textcolor{quarto-callout-important-color}{\faExclamation}
\end{minipage}%
\begin{minipage}[t]{\textwidth - 5.5mm}

\textbf{KVKK (Kişisel Verilerin Korunması Kanunu) / veri sınırı.} Bu
belge ham görüşme dökümlerini, doğrudan alıntı metinlerini, kimlikleyici
bilgileri veya aile-düzeyinde geri-tanımlanabilir ayrıntıları
\textbf{içermez}. Alıntı kanıtı yalnız \texttt{quote\_id} ile
gösterilmiştir. Tez metnine konacak doğrudan alıntılar, temizlenmiş tez
metninden \texttt{quote\_id} eşlemesiyle bire bir alınır, yeniden
yazılmaz. AI destekli çıktılar yalnız yardımcı ön-denetim/tutarlılık
kontrolüdür; kodlama, tema geliştirme ve yorumlama kararları araştırmacı
sorumluluğundadır.

\end{minipage}%
\end{tcolorbox}

\chapter{Yönetici Özeti}\label{sec-ozet}

Bu niteliksel çalışma, Tip 1 Diyabet Mellitus (T1DM) tanısı alan
çocuklar, sağlıklı kardeşleri ve anneleriyle yürütülen bireysel yarı
yapılandırılmış görüşmelere dayalı \textbf{triadik aile çalışması}dır.
Bulgular, T1DM'nin aile içinde tek ve homojen bir deneyim olarak
yaşanmadığını göstermektedir. Aynı tanı, anne için tıbbi bakım, risk
yönetimi, suçluluk ve sürekli tetikte olma sorumluluğuna; T1DM tanılı
çocuk için bedensel rutinler, sosyal görünürlük, normalleşme ve özerklik
müzakeresine; sağlıklı kardeş için görünmeyen kısıtlanma, erken
sorumluluk, adalet gerilimi ve sessiz feragat deneyimine dönüşmektedir.

Çalışmanın ana bilimsel sonucu şudur: T1DM, aile içi rolleri ve gündelik
yaşamı yalnız hasta çocuk etrafında değil, tüm aile üyelerinin bakım,
kontrol, adalet ve görünürlük pozisyonlarını yeniden kurarak
şekillendirmektedir. Anne, hasta çocuk ve sağlıklı kardeş aynı aile
olayını farklı risk, görev ve hakkaniyet çerçeveleriyle
anlamlandırmaktadır. Bu nedenle niteliksel sonuçlar, nicel tez
bulgularını ``doğrulamak'' için değil, anne--çocuk--kardeş algı
ayrışmalarını ve ölçeklerin yakalayamayabileceği deneyim katmanlarını
bağlamsallaştırmak için kullanılmalıdır.

Kanonik bulgu yapısı dört tez makro temasından oluşur:

\begin{enumerate}
\def\labelenumi{\arabic{enumi}.}
\tightlist
\item
  Sağlıklı kardeşin görünmeyen yükü.
\item
  Annenin tıbbi bakıcı rolüne kayması.
\item
  T1DM tanılı çocuğun içeriden hastalık deneyimi.
\item
  Aynı evde üç farklı deneyim olarak triadik karşılaştırma.
\end{enumerate}

Bu dört tema, journal yayını için geliştirilmiş altı konu-temalı yapı
ile aynı codebook (kod kitabı) tabanına dayanır. Tez ve journal tema
yapıları birbirinin alternatifi değil, aynı verinin iki farklı bilimsel
raporlama kesitidir.

\chapter{Kanonik Kapsam ve Veri Durumu}\label{sec-kapsam}

\section{Çalışma Kimliği}\label{uxe7alux131ux15fma-kimliux11fi}

{\def\LTcaptype{none} % do not increment counter
\begin{longtable}[]{@{}
  >{\raggedright\arraybackslash}p{(\linewidth - 2\tabcolsep) * \real{0.5000}}
  >{\raggedright\arraybackslash}p{(\linewidth - 2\tabcolsep) * \real{0.5000}}@{}}
\toprule\noalign{}
\begin{minipage}[b]{\linewidth}\raggedright
Bileşen
\end{minipage} & \begin{minipage}[b]{\linewidth}\raggedright
Kanonik durum
\end{minipage} \\
\midrule\noalign{}
\endhead
\bottomrule\noalign{}
\endlastfoot
Çalışma türü & Niteliksel triadik aile çalışması \\
Klinik bağlam & T1DM tanılı çocuk, anne ve sağlıklı kardeş aile
deneyimi \\
Tasarım & Niteliksel tanımlayıcı yaklaşım + fenomenolojik duyarlılık +
multi-informant family design \\
Analiz yaklaşımı & Braun ve Clarke çizgisinde refleksif tematik analiz
(Braun ve Clarke, 2006; Braun ve Clarke, 2019) \\
Raporlama standardı & COREQ (Consolidated Criteria for Reporting
Qualitative Research) 32 madde uyumlu (Tong ve ark., 2007); SRQR
(Standards for Reporting Qualitative Research) (O'Brien ve ark., 2014)
ve JARS-Qual (Journal Article Reporting Standards for Qualitative
Research) (Levitt ve ark., 2018) ek denetim \\
Analiz birimi & Aile triadı ve rol temelli anlatı \\
Görüşme tipi & Bireysel yarı yapılandırılmış görüşme \\
Görüşme rolleri & Anne, T1DM tanılı çocuk, sağlıklı kardeş \\
Veri toplama bağlamı & Klinik birim içinde mahremiyeti sağlanmış görüşme
ortamı \\
Veri yönetimi & KVKK uyumlu yerel ve anonimleştirilmiş çalışma hattı \\
\end{longtable}
}

\section{Kanonik Sayılar}\label{kanonik-sayux131lar}

{\def\LTcaptype{none} % do not increment counter
\begin{longtable}[]{@{}lr@{}}
\toprule\noalign{}
Bileşen & Sayı \\
\midrule\noalign{}
\endhead
\bottomrule\noalign{}
\endlastfoot
Aile triadı & 7 \\
Toplam katılımcı & 21 \\
Anne & 7 \\
T1DM tanılı çocuk & 7 \\
Sağlıklı kardeş & 7 \\
Aile kodları & 011, 014, 019, 020, 026, 201, 202 \\
Tez makro teması & 4 \\
Journal teması & 6 \\
Codebook kodu & 23 \\
Seçilmiş illüstratif alıntı (quote ID) havuzu & 116 \\
Araştırmacı-denetimli kodlanmış segment & 116 \\
Triadik matris satırı (aile × tema × alt-tema) & 57 \\
COREQ maddesi & 32 \\
COREQ tam karşılanan & 30 \\
COREQ kısmi & 2 \\
COREQ eksik & 0 \\
\end{longtable}
}

\section{Rol ve Aile
Dağılımı}\label{rol-ve-aile-daux11fux131lux131mux131}

Seçilmiş quote ID havuzu, nihai kodlama yoğunluğunu değil, tez yazımında
kullanılacak kanıt havuzunun kapsamını gösterir.

{\def\LTcaptype{none} % do not increment counter
\begin{longtable}[]{@{}lr@{}}
\toprule\noalign{}
Rol & Seçilmiş quote sayısı \\
\midrule\noalign{}
\endhead
\bottomrule\noalign{}
\endlastfoot
Anne & 55 \\
T1DM tanılı çocuk & 36 \\
Sağlıklı kardeş & 25 \\
\end{longtable}
}

{\def\LTcaptype{none} % do not increment counter
\begin{longtable}[]{@{}lr@{}}
\toprule\noalign{}
Aile & Seçilmiş quote sayısı \\
\midrule\noalign{}
\endhead
\bottomrule\noalign{}
\endlastfoot
011 & 24 \\
014 & 18 \\
019 & 15 \\
020 & 14 \\
026 & 23 \\
201 & 8 \\
202 & 14 \\
\end{longtable}
}

{\def\LTcaptype{none} % do not increment counter
\begin{longtable}[]{@{}lrl@{}}
\toprule\noalign{}
Tez makro teması & Quote sayısı & Rol dağılımı \\
\midrule\noalign{}
\endhead
\bottomrule\noalign{}
\endlastfoot
Tema 1 & 3 & sağlıklı kardeş: 3 \\
Tema 2 & 40 & anne: 40 \\
Tema 3 & 21 & T1DM tanılı çocuk: 21 \\
Tema 4 & 52 & anne: 15; T1DM tanılı çocuk: 15; sağlıklı kardeş: 22 \\
\end{longtable}
}

Bu dağılım, Tema 1--3'ün bilinçli olarak perspektif-organize olduğunu,
Tema 4'ün ise triadik çapraz okuma için kullanıldığını gösterir. Tema
1--3'te tek rol baskınlığı metodolojik kusur değil, tez mimarisinin
sonucudur.

\chapter{Metodoloji}\label{sec-metodoloji}

\section{Araştırma Tasarımı ve Epistemolojik
Konum}\label{araux15ftux131rma-tasarux131mux131-ve-epistemolojik-konum}

Çalışma, aynı aile bağlamında üç bilgi verici pozisyonunu ayrı ayrı
görünür kılan \textbf{multi-informant family design} üzerine kuruludur.
Çoklu bilgi verici yaklaşımının geçerliliği ve bilgi vericiler arası
tutarsızlıkların ölçüm hatası değil bağlama özgü anlamlı bilgi olduğu
ilgili literatürde gösterilmiştir (De Los Reyes ve ark., 2015). Her
aile, anne, T1DM tanılı çocuk ve sağlıklı kardeşten oluşan bir bilgi
üçlüsü olarak ele alınmıştır. Görüşmeler birlikte değil, ayrı ayrı
bireysel görüşme biçiminde yapılmıştır. Ayrı bireysel görüşme tercihi,
aile içi güç asimetrisi, uyum baskısı ve çocuk/kardeş sesinin
gölgelenmesi riskini azaltmak için metodolojik olarak uygundur; birlikte
ve ayrı görüşme biçimlerinin yöntem otantikliği açısından farklı
sonuçlar doğurabildiği bilinmektedir (Taylor ve Vocht, 2011).

\begin{tcolorbox}[enhanced jigsaw, arc=.35mm, bottomrule=.15mm, bottomtitle=1mm, breakable, colback=white, colbacktitle=quarto-callout-note-color!10!white, colframe=quarto-callout-note-color-frame, coltitle=black, left=2mm, leftrule=.75mm, opacityback=0, opacitybacktitle=0.6, rightrule=.15mm, title=\textcolor{quarto-callout-note-color}{\faInfo}\hspace{0.5em}{Not}, titlerule=0mm, toprule=.15mm, toptitle=1mm]

\textbf{Bilgi kutusu --- Triadik multi-informant tasarım ve neden
\emph{ayrı} bireysel görüşme?}

\textbf{Multi-informant (çoklu bilgi verici) tasarım}, aynı aile olayını
farklı konumlardan gören bilgi vericileri birlikte inceler; bilgi
vericiler arası tutarsızlık ölçüm hatası değil, bağlama özgü anlamlı
bilgidir (De Los Reyes ve ark., 2015). Bu çalışmada üçlü (triad) =
\textbf{anne + T1DM tanılı çocuk + sağlıklı kardeş}. Görüşmeler birlikte
(conjoint) değil \textbf{ayrı ayrı bireysel} yapılmıştır: ayrı görüşme,
aile içi güç asimetrisini ve uyum baskısını azaltarak çocuğun/kardeşin
sesinin gölgelenmesini önler; birlikte ve ayrı görüşme biçimleri yöntem
otantikliği açısından farklı veri üretir (Taylor ve Vocht, 2011).

\end{tcolorbox}

Analiz, Braun ve Clarke'ın \textbf{refleksif tematik analiz (RTA)}
yaklaşımıyla yürütülmüştür (Braun ve Clarke, 2006; Braun ve Clarke,
2019). RTA'nın epistemolojik konumu üç sonucu zorunlu kılar: (i) tema,
veride keşfedilmeyi bekleyen sabit bir nesne değil, araştırmacı ile veri
arasında geliştirilen \textbf{merkezi düzenleyici kavramdır}; (ii)
araştırmacı anlam üretiminde aktif faildir, dolayısıyla refleksivite bir
tehdit değil, analitik kaynaktır; (iii) örneklem yeterliliği sayısal
doygunlukla değil \textbf{bilgi gücüyle} gerekçelenir (Malterud ve ark.,
2016). Bu gerekçelerle, kodlama-güvenirlik (coding reliability)
yaklaşımına ait iki-kodlayıcı uzlaşım/kappa prosedürü bu çalışmada
\textbf{amaçlanmamıştır} (bkz. Bölüm~\ref{sec-raporlama}).

Tasarımın bilimsel gerekçesi:

\begin{itemize}
\tightlist
\item
  T1DM yalnız hasta çocuğun biyomedikal durumu değildir; aile
  rutinlerini, ebeveynlik pratiğini, kardeş ilişkilerini ve risk
  yönetimini etkileyen aile-düzeyi bir deneyimdir; kronik hastalığa aile
  yanıtını durumun tanımı, yönetim davranışları ve algılanan sonuçlar
  ekseninde çözümleyen Aile Yönetim Tarzı Çerçevesi (FMSF) bu
  aile-düzeyi bakışın kuramsal dayanağıdır (Knafl ve Deatrick, 2003).
\item
  Anne, hasta çocuk ve sağlıklı kardeş aynı olayın farklı tanıklarıdır;
  bu fark tek bir ``doğru'' anlatıya indirgenmemelidir.
\item
  Triadik analiz, yalnız uyumlu anlatıları değil rol-temelli
  ayrışmaları, sessizleşen yükleri ve aile içi kör alanları görünür
  kılar.
\item
  Ayrı bireysel görüşme, her rolün kendi deneyim dilini korumayı
  amaçlar.
\end{itemize}

\section{Altı-Faz RTA Protokolü}\label{altux131-faz-rta-protokoluxfc}

\begin{tcolorbox}[enhanced jigsaw, arc=.35mm, bottomrule=.15mm, bottomtitle=1mm, breakable, colback=white, colbacktitle=quarto-callout-note-color!10!white, colframe=quarto-callout-note-color-frame, coltitle=black, left=2mm, leftrule=.75mm, opacityback=0, opacitybacktitle=0.6, rightrule=.15mm, title=\textcolor{quarto-callout-note-color}{\faInfo}\hspace{0.5em}{Not}, titlerule=0mm, toprule=.15mm, toptitle=1mm]

\textbf{Bilgi kutusu --- Refleksif tematik analiz (RTA) nedir?}

Refleksif tematik analiz, Braun ve Clarke'ın geliştirdiği, nitel veride
anlam örüntülerini \textbf{tema} olarak kavramsallaştıran yorumlayıcı
(Big-Q) bir yaklaşımdır (Braun ve Clarke, 2006; Braun ve Clarke, 2019).
Üç ayırt edici ilkesi vardır:

\begin{itemize}
\tightlist
\item
  \textbf{Tema = merkezi düzenleyici kavram.} Tema, veride
  ``keşfedilmeyi bekleyen'' hazır bir nesne değil; araştırmacının
  veriyle etkileşiminde \emph{geliştirdiği} bir anlam örüntüsüdür.
\item
  \textbf{Refleksivite bir kaynaktır, tehdit değil.} Araştırmacının
  öznelliği elenmesi gereken bir yanlılık değil, analizin derinliğini
  besleyen analitik bir araçtır.
\item
  \textbf{Altı faz özyinelemelidir, doğrusal değil.} Fazlar arasında
  ileri-geri hareket edilir; aşağıdaki tablo her fazın bu koldaki somut
  uygulamasını verir.
\end{itemize}

\end{tcolorbox}

{\def\LTcaptype{none} % do not increment counter
\begin{longtable}[]{@{}
  >{\raggedright\arraybackslash}p{(\linewidth - 4\tabcolsep) * \real{0.3333}}
  >{\raggedright\arraybackslash}p{(\linewidth - 4\tabcolsep) * \real{0.3333}}
  >{\raggedright\arraybackslash}p{(\linewidth - 4\tabcolsep) * \real{0.3333}}@{}}
\toprule\noalign{}
\begin{minipage}[b]{\linewidth}\raggedright
Faz
\end{minipage} & \begin{minipage}[b]{\linewidth}\raggedright
İş
\end{minipage} & \begin{minipage}[b]{\linewidth}\raggedright
Bu koldaki uygulama
\end{minipage} \\
\midrule\noalign{}
\endhead
\bottomrule\noalign{}
\endlastfoot
1. Aşinalık & Tam, tekrarlı okuma + refleksif not & Anonimleştirilmiş
çalışma metni + paralinguistik saha notu birlikte okundu \\
2. Kodlama & Veriye yakın açık kodlar & 116 araştırmacı-denetimli
kodlanmış segment \\
3. Tema oluşturma & Kodların kavramsal kümelenmesi & 23 kod → 5 kategori
→ 6 journal / 4 tez teması \\
4. Gözden geçirme & Temaları veriye/birbirine karşı test & Aile-içi +
aileler-arası okuma ile tutarlılık sınandı \\
5. Tanımlama + isimlendirme & Her temaya kavramsal kimlik & Tema 1--4
makro kimlikleri \\
6. Yazım & Analitik anlatı + \texttt{quote\_id} ankraj &
Bölüm~\ref{sec-bulgular} ve bu belge \\
\end{longtable}
}

\section{Örnekleme ve Bilgi
Gücü}\label{uxf6rnekleme-ve-bilgi-guxfccuxfc}

Örneklem, sayısal doygunluk eşiği yerine \textbf{bilgi gücü}
çerçevesiyle gerekçelenmiştir; bu tercih RTA epistemolojisiyle uyumludur
(Malterud ve ark., 2016; Braun ve Clarke, 2019).

\begin{tcolorbox}[enhanced jigsaw, arc=.35mm, bottomrule=.15mm, bottomtitle=1mm, breakable, colback=white, colbacktitle=quarto-callout-note-color!10!white, colframe=quarto-callout-note-color-frame, coltitle=black, left=2mm, leftrule=.75mm, opacityback=0, opacitybacktitle=0.6, rightrule=.15mm, title=\textcolor{quarto-callout-note-color}{\faInfo}\hspace{0.5em}{Not}, titlerule=0mm, toprule=.15mm, toptitle=1mm]

\textbf{Bilgi kutusu --- ``Bilgi gücü'' (information power) neden
doygunluğun yerini alır?}

Malterud ve arkadaşlarının önerdiği \textbf{bilgi gücü}, bir nitel
örneklemin yeterliliğini sayısal bir eşikle değil, verinin araştırma
sorusuna yanıt üretme \emph{kapasitesiyle} değerlendirir (Malterud ve
ark., 2016). Örneklem ne kadar çok bilgi gücü taşırsa gereken katılımcı
sayısı o kadar azalır. Beş boyut birlikte tartılır: (1) çalışmanın
\textbf{amaç darlığı}, (2) örneklem \textbf{özgüllüğü}, (3) yerleşik
\textbf{kuram} kullanımı, (4) \textbf{diyalog kalitesi}, (5)
\textbf{analiz stratejisi}. Bu çalışma dört boyutta yüksek, birinde
orta-yüksek profil gösterdiğinden yedi aile triadı bilgi açısından
yeterli kabul edilmiştir. ``Doygunluk'' dili bilinçli olarak
kullanılmamıştır: refleksif TA'nın yorumsal duruşunda tema
\emph{üretilir}, tüketilene dek \emph{keşfedilmez} (Braun ve Clarke,
2021).

\end{tcolorbox}

{\def\LTcaptype{none} % do not increment counter
\begin{longtable}[]{@{}
  >{\raggedright\arraybackslash}p{(\linewidth - 4\tabcolsep) * \real{0.3333}}
  >{\raggedright\arraybackslash}p{(\linewidth - 4\tabcolsep) * \real{0.3333}}
  >{\raggedright\arraybackslash}p{(\linewidth - 4\tabcolsep) * \real{0.3333}}@{}}
\toprule\noalign{}
\begin{minipage}[b]{\linewidth}\raggedright
Bilgi gücü boyutu
\end{minipage} & \begin{minipage}[b]{\linewidth}\raggedright
Bu çalışmadaki durum
\end{minipage} & \begin{minipage}[b]{\linewidth}\raggedright
Değerlendirme
\end{minipage} \\
\midrule\noalign{}
\endhead
\bottomrule\noalign{}
\endlastfoot
Amaç darlığı & T1DM ailelerinde anne--hasta--kardeş deneyimine odak;
genel kronik hastalık değil & Yüksek \\
Örneklem özgüllüğü & Tüm ailelerde T1DM'li çocuk + sağlıklı kardeş; tek
klinik amaçlı örneklem & Yüksek \\
Kuramsal çerçeve & Multi-informant family design (De Los Reyes ve ark.,
2015) + pediatrik kronik hastalık aile yükü (Knafl ve Deatrick, 2003) &
Orta-yüksek \\
Görüşme kalitesi & Klinik deneyimli araştırmacı, mahrem ortam, bireysel
görüşme, paralinguistik not & Yüksek \\
Analiz stratejisi & Aile-içi triadik + aileler-arası tema
karşılaştırması & Orta-yüksek \\
\end{longtable}
}

\textbf{Sonuç:} Yedi aile triadı ve 21 bireysel görüşme, dar amaç, özgül
örneklem ve triadik analiz stratejisi nedeniyle bilgi açısından yoğun
bir veri seti olarak değerlendirilmiştir. Bu ifade istatistiksel
genelleme iddiası taşımaz. ``Doygunluğa ulaşıldı'' dili \textbf{bu
çalışmada} tercih edilmemiştir; çünkü refleksif TA'nın yorumsal (Big-Q)
epistemolojisinde tema araştırmacı ile veri arasında üretilir,
``tüketilene kadar keşfedilecek'' sabit bir nesne değildir (Braun ve
Clarke, 2019). Braun ve Clarke, doygunluk kavramının refleksif TA için
epistemolojik olarak tutarsız olduğunu ayrıca doğrudan tartışır (Braun
ve Clarke, 2021). Doygunluğun tematik analiz iş akışı içinde
operasyonelleştirilebileceğini savunan bir karşı-literatür bulunduğu
kabul edilir (Saunders ve ark., 2018); bu tezde bilgi gücü tercihi,
doygunluğun ``yanlış'' olduğu iddiasıyla değil, \textbf{refleksif TA ile
paradigma tutarlılığı} gerekçesiyle yapılmıştır.

\section{Veri Toplama}\label{veri-toplama}

{\def\LTcaptype{none} % do not increment counter
\begin{longtable}[]{@{}
  >{\raggedright\arraybackslash}p{(\linewidth - 2\tabcolsep) * \real{0.5000}}
  >{\raggedright\arraybackslash}p{(\linewidth - 2\tabcolsep) * \real{0.5000}}@{}}
\toprule\noalign{}
\begin{minipage}[b]{\linewidth}\raggedright
Özellik
\end{minipage} & \begin{minipage}[b]{\linewidth}\raggedright
Kanonik raporlama
\end{minipage} \\
\midrule\noalign{}
\endhead
\bottomrule\noalign{}
\endlastfoot
Görüşme formu & Anne, T1DM'li çocuk, sağlıklı kardeş için yarı
yapılandırılmış rehber \\
Görüşme sıklığı & Her katılımcı ile tek görüşme \\
Tekrarlı görüşme & Planlanmamış \\
Görüşme süresi & Bireysel \textasciitilde15--25 dakika; aile toplamı
\textasciitilde40--70 dakika \\
Ortam & Çocuk endokrinoloji birimine ait sessiz ve mahrem görüşme
odası \\
Görüşmede bulunanlar & Katılımcı, OM ve BA; başka aile üyesi/personel
yok \\
Kayıt & Ses kaydı + eşzamanlı saha/paralinguistik not \\
Transkript iadesi & Formal iade yok; görüşme-içi özetleme ile katılımcı
doğrulaması \\
\end{longtable}
}

Ayrı bir formal pilot görüşme yürütülmemiştir; yarı-yapılandırılmış
rehber alanyazın ve klinik deneyimle geliştirilip saha sürecinde
iyileştirilmiştir (COREQ Madde 16; bkz. Bölüm~\ref{sec-raporlama}).

\section{Araştırmacı Rolleri ve
Refleksivite}\label{araux15ftux131rmacux131-rolleri-ve-refleksivite}

OM birinci araştırmacı ve temel kodlayıcıdır. BA görüşme sırasında
non-participant observer olarak paralinguistik not tutmuş, analizde
\textbf{critical friend} rolü üstlenmiştir. RTA epistemolojisi gereği
inter-coder agreement ya da kappa hesabı amaçlanmamıştır; BA'nın rolü
kodları ``doğrulamak'' değil, alternatif okuma ve refleksif sorgulama
sağlamaktır (Braun ve Clarke, 2019).

{\def\LTcaptype{none} % do not increment counter
\begin{longtable}[]{@{}
  >{\raggedright\arraybackslash}p{(\linewidth - 2\tabcolsep) * \real{0.5000}}
  >{\raggedright\arraybackslash}p{(\linewidth - 2\tabcolsep) * \real{0.5000}}@{}}
\toprule\noalign{}
\begin{minipage}[b]{\linewidth}\raggedright
Refleksif risk
\end{minipage} & \begin{minipage}[b]{\linewidth}\raggedright
Yönetim stratejisi
\end{minipage} \\
\midrule\noalign{}
\endhead
\bottomrule\noalign{}
\endlastfoot
Hekim kimliğinin güç asimetrisi & Klinik tedavi ilişkisi olmaması, rol
açıklaması, önlüksüz görüşme, mahrem ortam \\
Klinik aşırı-yorum & Katılımcı diline yakın kodlama, reflexive memo,
critical friend tartışması \\
Empatik aşırı-ağırlıklandırma & Paralinguistik notların açık kullanımı,
alternatif okuma \\
Disiplin-içi ortak kör nokta & BA critical friend, klinik-dışı okuma
denemeleri \\
Gözlemci varlığının etkisi & BA rolünün görüşme öncesi açıklanması, arka
planda kalma \\
\end{longtable}
}

\section{Analiz Yaklaşımı ve Araç
Sınırı}\label{analiz-yaklaux15fux131mux131-ve-arauxe7-sux131nux131rux131}

Kodlama ve tema geliştirme manuel araştırmacı yorumu temelindedir. Yerel
toolkit (\texttt{dmnitel}) tarafından üretilen CSV ve raporlar
\textbf{nihai analiz üretmez}; yalnız denetim, düzenleme ve tutarlılık
kontrolü sağlar. Analiz zinciri: (1) görüşmelerin dökülmesi ve
anonimleştirilmiş çalışma metnine dönüştürülmesi; (2) veriye yakın açık
kodlar; (3) 5 kategori ve 23 kodluk codebook düzeni; (4) journal 6-tema
ve tez 4-makro tema yapılarının aynı codebook tabanına bağlanması; (5)
aile-içi triadik matrislerle anne--hasta çocuk--sağlıklı kardeş
anlatılarının karşılaştırılması; (6) aileler-arası karşılaştırma; (7)
critical friend tartışmaları + audit trail; (8) quote integrity ve
negatif vaka analizi ile kalite kontrol.

\chapter{Etik, KVKK ve Veri Güvenliği}\label{sec-etik}

Çalışma özel nitelikli \textbf{sağlık + çocuk} verisi içerdiğinden
yüksek gizlilik düzeyiyle yürütülmüştür. Ham ses kayıtları,
anonimleştirilmemiş transkriptler, demografik formlar, onam belgeleri,
pseudonym haritaları ve aile-düzeyinde geri-tanımlanabilir satırlar
rapora, OSF'ye veya büyük dil modeli (LLM) sistemlerine aktarılmaz.

{\def\LTcaptype{none} % do not increment counter
\begin{longtable}[]{@{}
  >{\raggedright\arraybackslash}p{(\linewidth - 2\tabcolsep) * \real{0.5000}}
  >{\raggedright\arraybackslash}p{(\linewidth - 2\tabcolsep) * \real{0.5000}}@{}}
\toprule\noalign{}
\begin{minipage}[b]{\linewidth}\raggedright
Alan
\end{minipage} & \begin{minipage}[b]{\linewidth}\raggedright
İlke
\end{minipage} \\
\midrule\noalign{}
\endhead
\bottomrule\noalign{}
\endlastfoot
Ham veri & Yerel, şifreli, erişimi sınırlı saklanır \\
Anonimleştirme & İsim, adres, doğum tarihi ve doğrudan tanımlayıcılar
kaldırılır/maskelenir \\
Pseudonym haritası & Ayrı şifreli dosyada; tez sonrası belirlenen sürede
silinir \\
Yayın alıntıları & Rol etiketi + aile kodu ile, k-anonymity
gözetilerek \\
Coğrafi/demografik bilgi & Bantlama + toplulaştırma ile raporlanır \\
Çocuk assent & Ebeveyn onamı + gelişim düzeyine uygun çocuk assent;
katılımın geri çekilebilirliği \\
OSF & Ham veri yüklenmez; codebook, metodoloji ve anonim raporlama
materyali paylaşılabilir \\
LLM kullanımı & Yazım/denetim desteği; ham veya kimliklenebilir veri
gönderimi yok \\
\end{longtable}
}

\textbf{AI/LLM kullanım beyanı.} LLM kodlama, tema geliştirme veya nihai
yorumlama yapmamıştır; yazım düzenleme, metin organizasyonu, raporlama
uyumu ve ön-denetim desteği için kullanılmıştır. Ham transcript, saha
notu, demografik form, pseudonym haritası ve tanımlayıcı bilgi LLM'e
gönderilmez. LLM çıktıları OM tarafından gözden geçirilmeden akademik
metne aktarılmaz (bkz. Bölüm~\ref{sec-yonetisim}).

\chapter{Codebook ve Analitik Kavram Haritası}\label{sec-codebook}

Codebook 23 koddan oluşur. Bu kodların 18'i seçilmiş illüstratif alıntı
havuzunda ankrajlıdır; kalan beş kod bu havuzda temsil edilmez:
\texttt{TANI\_BILGI\_ARAYISI}, \texttt{GECE\_TAKIP},
\texttt{YARAR\_BULMA\_OLGUNLASMA}, \texttt{SAGLIK\_EKIBI\_DESTEGI},
\texttt{MANEVI\_BASA\_CIKMA}. Bu beş kod codebook'un kavramsal yapısında
yer alır; illüstratif havuz tüm kodları tüketmeyi değil, dört makro
temanın anlatısını desteklemeyi amaçlayan purposive bir seçkidir (bkz.
Bölüm~\ref{sec-sinirliliklar}).

\section{Kod Kategorileri}\label{kod-kategorileri}

{\def\LTcaptype{none} % do not increment counter
\begin{longtable}[]{@{}
  >{\raggedright\arraybackslash}p{(\linewidth - 2\tabcolsep) * \real{0.5000}}
  >{\raggedright\arraybackslash}p{(\linewidth - 2\tabcolsep) * \real{0.5000}}@{}}
\toprule\noalign{}
\begin{minipage}[b]{\linewidth}\raggedright
Kategori
\end{minipage} & \begin{minipage}[b]{\linewidth}\raggedright
Kodlar
\end{minipage} \\
\midrule\noalign{}
\endhead
\bottomrule\noalign{}
\endlastfoot
Tanı ve ilk dönem deneyimi & \texttt{TANI\_SOK\_KORKU},
\texttt{TANI\_HASTANE\_SURECI}, \texttt{TANI\_BILGI\_ARAYISI} \\
Günlük yönetim yükü ve rutinlerin yeniden örgütlenmesi &
\texttt{RUTIN\_TAKIP}, \texttt{BESLENME\_KONTROL}, \texttt{GECE\_TAKIP},
\texttt{OKUL\_SOSYAL\_UYUM}, \texttt{TEKNOLOJI\_DESTEGI} \\
Duygusal deneyim ve anlamlandırma & \texttt{KAYGI\_KIRILGANLIK},
\texttt{SUCLULUK}, \texttt{OFKE\_ADALETSIZLIK},
\texttt{KABULLENME\_NORMALLESME}, \texttt{YARAR\_BULMA\_OLGUNLASMA} \\
Aile ilişkileri ve rollerin dönüşümü & \texttt{ANNE\_HIPERVIJILANS},
\texttt{COCUK\_OZERKLIK\_OZBAKIM}, \texttt{KARDES\_GORUNMEZ\_YUK},
\texttt{KARDES\_ILISKISI}, \texttt{AILE\_ICI\_ADALET},
\texttt{ILETISIM\_CATISMA} \\
Başa çıkma kaynakları ve destek sistemleri & \texttt{AILE\_DESTEGI},
\texttt{SAGLIK\_EKIBI\_DESTEGI}, \texttt{AKRAN\_CEVRE\_DESTEGI},
\texttt{MANEVI\_BASA\_CIKMA} \\
\end{longtable}
}

\section{Kod Yoğunluğu (denetim
sayısı)}\label{kod-youx11funluux11fu-denetim-sayux131sux131}

Aşağıdaki sayılar nihai tematik önem sıralaması \textbf{değildir};
temizlenmiş tez metninden çıkarılan seçilmiş quote ID havuzunun denetim
sayılarıdır. Niteliksel önem, yalnız frekansla değil temanın araştırma
sorusundaki merkezi düzenleyici işleviyle değerlendirilir.

{\def\LTcaptype{none} % do not increment counter
\begin{longtable}[]{@{}lrlr@{}}
\toprule\noalign{}
Kod & Ön-kodlu segment & Kod & Ön-kodlu segment \\
\midrule\noalign{}
\endhead
\bottomrule\noalign{}
\endlastfoot
\texttt{KARDES\_GORUNMEZ\_YUK} & 21 & \texttt{AILE\_DESTEGI} & 5 \\
\texttt{RUTIN\_TAKIP} & 16 & \texttt{KABULLENME\_NORMALLESME} & 5 \\
\texttt{KAYGI\_KIRILGANLIK} & 11 & \texttt{AILE\_ICI\_ADALET} & 4 \\
\texttt{OFKE\_ADALETSIZLIK} & 11 & \texttt{SUCLULUK} & 3 \\
\texttt{TANI\_SOK\_KORKU} & 10 & \texttt{ILETISIM\_CATISMA} & 2 \\
\texttt{KARDES\_ILISKISI} & 8 & \texttt{TEKNOLOJI\_DESTEGI} & 2 \\
\texttt{COCUK\_OZERKLIK\_OZBAKIM} & 7 & \texttt{OKUL\_SOSYAL\_UYUM} &
1 \\
\texttt{BESLENME\_KONTROL} & 7 & \texttt{AKRAN\_CEVRE\_DESTEGI} & 1 \\
& & \texttt{TANI\_HASTANE\_SURECI} & 1 \\
& & \texttt{ANNE\_HIPERVIJILANS} & 1 \\
\end{longtable}
}

\chapter{Tema Mimarisi}\label{sec-tema-mimarisi}

\section{Tez İçin 4 Makro Tema}\label{tez-iuxe7in-4-makro-tema}

{\def\LTcaptype{none} % do not increment counter
\begin{longtable}[]{@{}
  >{\raggedright\arraybackslash}p{(\linewidth - 6\tabcolsep) * \real{0.2500}}
  >{\raggedright\arraybackslash}p{(\linewidth - 6\tabcolsep) * \real{0.2500}}
  >{\raggedright\arraybackslash}p{(\linewidth - 6\tabcolsep) * \real{0.2500}}
  >{\raggedright\arraybackslash}p{(\linewidth - 6\tabcolsep) * \real{0.2500}}@{}}
\toprule\noalign{}
\begin{minipage}[b]{\linewidth}\raggedright
Tema
\end{minipage} & \begin{minipage}[b]{\linewidth}\raggedright
Analitik odak
\end{minipage} & \begin{minipage}[b]{\linewidth}\raggedright
Alt yapı
\end{minipage} & \begin{minipage}[b]{\linewidth}\raggedright
Tez işlevi
\end{minipage} \\
\midrule\noalign{}
\endhead
\bottomrule\noalign{}
\endlastfoot
Tema 1. Gölgede Kalan Çocuklar: Paylaşılan Kısıtlılık ve Görünmez Yük &
Sağlıklı kardeş deneyimi & Gönüllü mahrumiyet; erken olgunlaşma/nöbetçi
kardeş; adalet ve baskı boyutları & Kardeş yükünü ve görünmeyen bakım
etkisini açar \\
Tema 2. Annelikte Eksen Kayması: Sağlıklı Çocuk Annesinden Tıbbi
Bakıcıya & Anne deneyimi & Suçluluk; tıbbi bakıcı rolü; kaybetme
korkusu; özerklik çatışması; eş rolünün silikleşmesi; adalet ikilemi &
Anne öz-raporu ve bakım yükünü bağlamsallaştırır \\
Tema 3. Hastalığın İçinden: T1DM Tanılı Çocuğun Yaşam Deneyimi & Hasta
çocuk deneyimi & Normalleşme; tedavi yükü; damgalanma; korunma-kontrol
ve kardeşlik & Çocuk algısı ve özerklik gerilimini açıklar \\
Tema 4. Aynı Ev, Üç Farklı Deneyim & Triadik çapraz okuma & Diyabetik
düzen; tetikte olma; ilgi/adalet/kontrol; küçük bakıcılar ve hasta
otoritesi & Karma tezin ana nitel köprü temasını kurar \\
\end{longtable}
}

\section{Journal İçin 6 Konu
Teması}\label{journal-iuxe7in-6-konu-temasux131}

{\def\LTcaptype{none} % do not increment counter
\begin{longtable}[]{@{}
  >{\raggedright\arraybackslash}p{(\linewidth - 2\tabcolsep) * \real{0.5000}}
  >{\raggedright\arraybackslash}p{(\linewidth - 2\tabcolsep) * \real{0.5000}}@{}}
\toprule\noalign{}
\begin{minipage}[b]{\linewidth}\raggedright
Journal teması
\end{minipage} & \begin{minipage}[b]{\linewidth}\raggedright
Tezdeki ana karşılık
\end{minipage} \\
\midrule\noalign{}
\endhead
\bottomrule\noalign{}
\endlastfoot
J1. Beslenme ve gündelik yaşamın yeniden örgütlenmesi & Tema 1, Tema 2,
Tema 3 ve Tema 4.1 \\
J2. Duygulanım ve kaygı iklimi & Tema 2.1--2.3 ve Tema 4.2 \\
J3. Algılanan adalet, ilgi asimetrisi ve kardeşlik dinamikleri & Tema
1.3, Tema 2.6 ve Tema 4.3 \\
J4. Sosyal çevre, görünürlük ve mahremiyet & Tema 3.3, Tema 2.5 ve Tema
4 \\
J5. Destek ekosistemi ve yapısal belirleyiciler & Tema 2, Tema 3 ve
Tartışma \\
J6. Kontrol-özerklik müzakeresi & Tema 2.4, Tema 3.4 ve Tema 4.4 \\
\end{longtable}
}

\section{Odak Aile Matrisi}\label{odak-aile-matrisi}

{\def\LTcaptype{none} % do not increment counter
\begin{longtable}[]{@{}llr@{}}
\toprule\noalign{}
Aile & Birincil odak olduğu journal alt-temalar & Toplam \\
\midrule\noalign{}
\endhead
\bottomrule\noalign{}
\endlastfoot
011 & J1.2, J2.2, J3.3, J5.1, J6.4 & 5 \\
014 & J2.1, J3.1, J6.1 & 3 \\
019 & J1.1, J2.4, J4.1 & 3 \\
020 & J2.3, J4.2, J6.2 & 3 \\
026 & J3.4, J4.4, J5.3 & 3 \\
201 & J5.2 & 1 \\
202 & J1.3, J3.2, J4.3, J5.4, J6.3 & 5 \\
\end{longtable}
}

Aile 201'in düşük odak temsili, veri zenginliği ve aktarılabilirlik
tartışmasında açıkça ele alınmalıdır. Bu durum çalışmayı geçersiz
kılmaz; bilgi gücünün aileler arasında eşit dağılmadığını gösterir.

\chapter{Kanonik Bilimsel Bulgular}\label{sec-bulgular}

\section{Tema 1 --- Sağlıklı Kardeşin Görünmeyen Yükü}\label{tema-1}

\textbf{Ana sonuç.} Sağlıklı kardeş, T1DM tanısının doğrudan biyomedikal
öznesi olmasa da aile rutinlerindeki değişimi ve ebeveyn ilgisinin
yeniden dağılımını kendi gündelik yaşamında taşımaktadır. Kardeşin
deneyimi genellikle açık bir ``şikayet'' biçiminden çok, gönüllü
vazgeçme, kendini geri çekme, bakım sorumluluğunu normalleştirme ve
adalet gerilimini örtük taşıma şeklinde görünür.

{\def\LTcaptype{none} % do not increment counter
\begin{longtable}[]{@{}
  >{\raggedright\arraybackslash}p{(\linewidth - 4\tabcolsep) * \real{0.3333}}
  >{\raggedright\arraybackslash}p{(\linewidth - 4\tabcolsep) * \real{0.3333}}
  >{\raggedright\arraybackslash}p{(\linewidth - 4\tabcolsep) * \real{0.3333}}@{}}
\toprule\noalign{}
\begin{minipage}[b]{\linewidth}\raggedright
Alt tema
\end{minipage} & \begin{minipage}[b]{\linewidth}\raggedright
Bilimsel netice
\end{minipage} & \begin{minipage}[b]{\linewidth}\raggedright
Tezdeki işlev
\end{minipage} \\
\midrule\noalign{}
\endhead
\bottomrule\noalign{}
\endlastfoot
Gönüllü mahrumiyet ve vekaleten üzüntü & Kardeş, diyabetli kardeşin
kısıtlarını aile düzeninin ortak kuralı gibi yaşayabilir; dayanışma ile
kendi isteğini erteleme sınırında durur & Kardeş yükünü çatışma değil,
sessiz uyum ve feragat olarak göstermeye yarar \\
Erken olgunlaşma ve nöbetçi kardeş rolü & Erken
sorumluluk/izleme/hatırlatma; görünür kılınmadığında ``görünmeyen bakım
emeği''ne dönüşür & H2 kardeş ilişkisi bulgularında görünmeyen rol
yükü \\
İlgi asimetrisi ve ikinci planda kalma & Kabul, adalet duygusunun
zedelenmesini ortadan kaldırmaz & Anne niyeti ile kardeş deneyimi
farkı \\
Önleyici sağlık baskısı ve bedensel özerklik & Kardeş de
yiyecek/oyun/etkinlik açısından diyabetli kardeşin düzenine
uyarlanabilir & T1DM'nin hasta olmayan çocuğun beden/rutinine etkisi \\
Sosyal kıyaslanma ve kimlik yorgunluğu & Sürekli hasta kardeş üzerinden
konumlanma ikincil kimlik deneyimi yaratabilir & Kardeşin görünürlüğü \\
Hastalığın güç aracı olarak kullanımı & Hastalık kardeşler arası
güç/pazarlık ilişkisini etkileyebilir & Kardeşlik dinamiklerinin dengeli
okuması \\
\end{longtable}
}

\textbf{Bilimsel yorum.} Tema 1'in en önemli katkısı, sağlıklı kardeşin
``etkilenmeyen aile üyesi'' olarak ele alınamayacağını göstermesidir.
Nicel kardeş ilişkisi ölçümleri grup düzeyinde güçlü fark göstermese
bile, nitel bulgular kardeş yükünün ilişki kalitesi toplam puanından çok
bağlamsal durumlara, bakım rutinlerine, ebeveyn dikkatine ve adalet
algısına gömülü olabileceğini düşündürür.

Bu okuma dış literatürle güçlü biçimde örtüşür. Kronik hastalıklı
çocukların kardeşlerinde 51 çalışma ve 103 etki büyüklüğünü birleştiren
meta-analiz, kardeş uyumunda küçük ama olumsuz bir etki saptamış; kritik
olarak ebeveyn raporlarının çocukların öz-bildirimlerinden daha olumsuz
olduğunu (bilgi verici tutarsızlığının erken bir işareti) ve gündelik
tedavi rejimi gerektiren hastalıkların --- T1DM tipik bir örnektir ---
en olumsuz etkiyle ilişkilendiğini bildirmiştir (Sharpe ve Rossiter,
2002). T1DM'ye özgü güncel bir çalışma bu örüntüyü doğrudan
somutlaştırır: 100 kardeş ve 100 ebeveynden oluşan bir örneklemde
kardeşlerin \%38'i glisemik kontrol düzeyi, daha yüksek ebeveynlik stresi ve
kötü aile işlevselliğiyle ilişkili anlamlı duygusal-davranışsal sorunlar
bildirmiş; kardeş uyumu aile-bağıntılı bulunmuş ve uyumsal başa çıkma
bir tampon işlevi görmüştür (Elhabashy ve ark., 2023).

Kardeş yükünün toplam puandan çok adalet algısına gömülü olduğu
yorumumuz, algılanan ebeveyn-ilgisi asimetrisini doğrudan inceleyen
çalışmalarla desteklenir: kardeşler, ebeveynlerinin kendilerine hasta
kardeşlerinden daha az sevgi gösterdiğini algıladıklarında daha düşük
psikososyal ve fiziksel işlevsellik bildirmiştir (Kelada ve ark., 2022).
Daha eski kesitsel bir çalışma, T1DM'li çocuğun kardeşlerinin yarıdan
fazlasının (\%55) insülin ve diyet yönetimine yakından katıldığını, üçte
birinin (\%33) kendisini ailede en çok suçlanacak üye olarak gördüğünü
ve her iki durumun da düşük benlik saygısıyla ilişkili olduğunu
göstermiştir; tasarımın kesitsel niteliği nedeniyle nedensel bir yön
kurulamaz (Adams ve ark., 1991). Ebeveyn perspektifli bir çalışmada
aileler, kardeşlerin diyabet iş yükünü paylaştığını, hastalığın
kardeşlerde duygusal bir bedel oluşturduğunu ve diyabeti kardeşin
ihtiyaçlarının önüne koymaktan suçluluk duyduklarını dile getirmiş (Cao
ve ark., 2021); 23 nitel çalışmayı sentezleyen güncel bir derleme ise
sağlıklı kardeşlerin tanı sonrası psikolojik distres riski altında
olduğunu ve ``çok erken büyümek'' gibi temaların öne çıktığını ortaya
koymuştur (Tan ve ark., 2024). Bu kaynaklar deneyimi ``parentifikasyon''
veya ``algılanan adaletsizlik'' terimleriyle adlandırmaz;
bulgularımızdaki erken sorumluluk, nöbetçi kardeş rolü ve sessiz feragat
örüntüsünün bu kavramlarla okunması araştırmacının yorumsal katkısıdır.

\emph{Öne çıkan kodlar:} \texttt{KARDES\_GORUNMEZ\_YUK},
\texttt{OFKE\_ADALETSIZLIK}, \texttt{KARDES\_ILISKISI},
\texttt{AILE\_ICI\_ADALET}, \texttt{BESLENME\_KONTROL}. \emph{Quote ID
ankrajları:} \texttt{019\_healthy\_sibling\_q001},
\texttt{026\_healthy\_sibling\_q001},
\texttt{026\_healthy\_sibling\_q002}.

\section{Tema 2 --- Annenin Tıbbi Bakıcı Rolüne Kayması}\label{tema-2}

\textbf{Ana sonuç.} Tanı sonrası annelik, duygusal bakım veren
ebeveynlikten sürekli ölçüm, izlem, risk önleme, karar verme ve aile
düzenini yeniden kurma sorumluluğu taşıyan \textbf{tıbbi bakım
koordinatörlüğüne} kaymaktadır. Anne anlatılarında suçluluk, kaybetme
korkusu, tetikte olma, ergenlikte kontrol-özerklik çatışması ve sağlıklı
kardeşe adil davranma ikilemi birlikte görünür.

{\def\LTcaptype{none} % do not increment counter
\begin{longtable}[]{@{}
  >{\raggedright\arraybackslash}p{(\linewidth - 4\tabcolsep) * \real{0.3333}}
  >{\raggedright\arraybackslash}p{(\linewidth - 4\tabcolsep) * \real{0.3333}}
  >{\raggedright\arraybackslash}p{(\linewidth - 4\tabcolsep) * \real{0.3333}}@{}}
\toprule\noalign{}
\begin{minipage}[b]{\linewidth}\raggedright
Alt tema
\end{minipage} & \begin{minipage}[b]{\linewidth}\raggedright
Bilimsel netice
\end{minipage} & \begin{minipage}[b]{\linewidth}\raggedright
Tezdeki işlev
\end{minipage} \\
\midrule\noalign{}
\endhead
\bottomrule\noalign{}
\endlastfoot
Tanı sonrası suçluluk döngüsü ve sarsılan annelik kimliği & Suçluluk
yalnız geçmişe değil gelecekteki bakım sorumluluğuna da bağlanır & H3
anne öz-rapor ve H4 depresyon bağlantısı \\
Sürekli nöbet hali ve tıbbi bakıcıya evrilen annelik & Sıradan
ebeveynlikten ayrışan izleme/planlama/risk önleme rolü & Bakım yükünün
görünmeyen bilişsel/duygusal boyutu \\
Kaybetme korkusu ve en kötü senaryolar & Hipoglisemi/gece
riski/belirsizlik gelecek hayalini kırılganlaştırır & Kaygı ikliminin
aileye yayılımı \\
Ergenlikle kesişen yönetim ve otonomi çatışması & Bağımsızlaşma isteği
güvenlik-kontrol gerilimini artırır & H1/H5 anne-çocuk algı ayrışması
zemini \\
Mahremiyetin geri çekilişi ve eş rolünün silikleşmesi & Bakım yükü
evlilik/partnerlik alanını daraltabilir & Aile sistemindeki rol
kayması \\
Sağlıklı kardeşle ilişkide adalet ikilemi & Hasta çocuğu önceleme, adil
davranma çabasını etik-duygusal ikileme çevirir & Tema 1 ve Tema 4
köprüsü \\
\end{longtable}
}

\textbf{Bilimsel yorum.} Tema 2, anne öz-raporlarının neden her zaman
dramatik grup farkı göstermeyebileceğine dair nitel bağlam sağlar. Anne,
tıbbi bakıcı rolünü içselleştirip normal ebeveynlik görevi gibi
sunabilir; ölçeklerde görülen sınırlı fark yükün yokluğu değil, bakımın
normalleştirilmiş ve ahlaki sorumluluk olarak içselleştirilmiş
biçimidir.

Bu yorum, T1DM anne bakım yükü yazınıyla yakından örtüşür. On sekiz
nitel çalışmayı sentezleyen bir kapsam derlemesi, bakım yükünün
cinsiyetlenmiş olduğunu; gece glukoz izleme sorumluluğunu ağırlıkla
annelerin üstlendiğini, babalardan daha fazla uyku yoksunluğu
bildirdiğini ve kadınların istihdamlarını kısıtlama veya değiştirme
zorunluluğunu daha sık dile getirdiğini göstermiştir (Azimi ve ark.,
2024). Güncel bir nitel çalışma bu yükün fenomenolojisini derinleştirir:
kan şekerini sürekli izleme ve kesintisiz tetikte olma zorunluluğu
ebeveynin kişisel ve mesleki yaşamını bozmakta, yıllar içinde tükenme,
uykusuzluk ve bilişsel güçlük üreterek bakımı süreğen bir işe
dönüştürmektedir (ana tema: ``kendini düğüm düğüm bağlama'') (Carlsund
ve ark., 2025). Bir başka nitel çalışma, karmaşık ve titiz gündelik
tıbbi bakımın aile dengesini bozduğunu doğrudan temalaştırır (Haghighi
Moghadam ve ark., 2022); bir diğeri ise bakım yükünün artan yalnızlık
hissiyle birlikte gittiğini bildirir (Kobos ve ark., 2023).

Annenin öz-raporda her zaman belirmeyen suçluluk ve tetikte olma katmanı
da yazında izlenebilir: hipoglisemi korkusu ebeveynlerde yaygındır ve
özellikle annelerde belirgindir (Barnard ve ark., 2010); tanı sürecini
``gecikmiş'' bir yol olarak deneyimleyen, belirtileri başka durumlara
bağlayıp geç başvuran ebeveynler çözülmemiş suçluluk ve kendini suçlama
duyguları taşır (Rankin ve ark., 2014). Bu son çalışma tanı-öncesi
süreçle sınırlıdır; suçluluğun tanı sonrası tıbbi tetikte olma
davranışına dönüşmesi, bulgularımızın önerdiği ancak bu kaynağın
doğrudan sınamadığı yorumsal bir köprüdür.

\emph{Öne çıkan kodlar:} \texttt{TANI\_SOK\_KORKU}, \texttt{SUCLULUK},
\texttt{RUTIN\_TAKIP}, \texttt{KAYGI\_KIRILGANLIK},
\texttt{ANNE\_HIPERVIJILANS}, \texttt{COCUK\_OZERKLIK\_OZBAKIM},
\texttt{AILE\_DESTEGI}, \texttt{ILETISIM\_CATISMA}. \emph{Quote ID
ankrajları:} \texttt{014\_mother\_q001}, \texttt{014\_mother\_q002},
\texttt{019\_mother\_q001}, \texttt{020\_mother\_q001},
\texttt{026\_mother\_q001}.

\section{Tema 3 --- T1DM Tanılı Çocuğun İçeriden Deneyimi}\label{tema-3}

\textbf{Ana sonuç.} T1DM tanılı çocuklar hastalığı hem gündelik yaşamın
olağan parçası haline getirmeye çalışmakta hem de tedavi yükü, bedensel
müdahale, sosyal görünürlük, damgalanma ve aile içi kontrol üzerinden
somut bir sınır olarak deneyimlemektedir. Normalleşme dili yükün yokluğu
anlamına gelmez; çoğu zaman yükü sürdürülebilir kılma stratejisidir.
Kronik hastalıkta normalleştirme, aile yönetim tarzı yazınında
tanımlanmış bir yönetim/başa çıkma süreci olarak kavramsallaştırılır ve
yükün yokluğuyla eşitlenmemelidir (Robinson, 1993; Deatrick ve ark.,
1999; Knafl ve Deatrick, 2003); bu çalışmada kavram, aile-düzeyi
kökeninden çocuğun bireysel deneyimine taşınırken bu düzey farkı
gözetilmiştir.

{\def\LTcaptype{none} % do not increment counter
\begin{longtable}[]{@{}
  >{\raggedright\arraybackslash}p{(\linewidth - 4\tabcolsep) * \real{0.3333}}
  >{\raggedright\arraybackslash}p{(\linewidth - 4\tabcolsep) * \real{0.3333}}
  >{\raggedright\arraybackslash}p{(\linewidth - 4\tabcolsep) * \real{0.3333}}@{}}
\toprule\noalign{}
\begin{minipage}[b]{\linewidth}\raggedright
Alt tema
\end{minipage} & \begin{minipage}[b]{\linewidth}\raggedright
Bilimsel netice
\end{minipage} & \begin{minipage}[b]{\linewidth}\raggedright
Tezdeki işlev
\end{minipage} \\
\midrule\noalign{}
\endhead
\bottomrule\noalign{}
\endlastfoot
Normalleştirme ve avantaja çevirme & Olağanlaştırma bazen güçlü başa
çıkma, bazen performatif normallik & Çocuk anlatısındaki direnç-uyum
birlikteliği \\
Somut kısıtlılıklar ve tedavi yükü & Ölçüm/enjeksiyon/beslenme/planlama
gündelik beden pratiği & Biyomedikal yönetimin somut yükü \\
Utanç, farklılık ve sosyal damgalanma & Okul/akran çevresinde görünürlük
ve farklılık hissi & Sosyal uyum ve kimlik boyutu \\
Korunma, kontrol ve kardeşlik & Artan ilgi hem güven hem kontrol; özel
muamele ve adalet gerilimi & H5 algı tutarlılığı ve Tema 4 triadik
gerilimi \\
\end{longtable}
}

\textbf{Bilimsel yorum.} Tema 3, çocuğun hasta kimliği ile gündelik
çocuk/ergen kimliği arasında sürekli müzakere yaptığını gösterir.
``Normal'' söylemi hastalığın etkisizliği değil; bağımsızlık, akranlara
benzeme ve aile kontrolünü yönetme stratejisi olarak okunmalıdır.

Bu okuma, T1DM'li çocuk ve ergenlerin yaşanmış deneyimine odaklanan
nitel çalışmalarla desteklenir. Gömülü kuram temelli bir çalışma,
ergenin hastalık deneyimini benliğin diyabet kimliğiyle ``sürgün
edilmesinden'' (self-banished) hastalığın benlikle uzlaştırılmasına
uzanan ``yeniden kendim olmak'' sürecinde çerçeveler; bu süreçte çocuk
pasif bir hasta değil, hastalığın anlamını sosyal etkileşim içinde
yeniden kuran bir faildir (Kim, 2022). Kenya'da yürütülen bir çalışma
``normale ulaşma çabası'' (striving for normality) temasında bu
müzakerenin gündelik biçimlerini belgeler: kimi çocuk sınıfta rahatça
enjeksiyon yapabilirken, kimi glukometresinin görünür olmasının yol
açacağı zarardan çekinerek hastalığını sır olarak taşır; aynı çalışma
insülini kendi başına yönetme becerisinin gurur ve olgunlaşma kaynağı
olduğunu, ancak bu özerklik kazanımının ``acınma'' isteğinden kaçınma ve
tedaviyi diskret yürütme arzusuyla eşzamanlı işlediğini göstermiştir
(Palmer ve ark., 2022).

Normalleştirme çabasının arkasındaki damgalanma katmanı da nitel olarak
somutlaşır. Adölesanlarda stigma yalnız açık alay değildir; sempati ve
``ilgi'' kılığına bürünen örtük önyargı biçiminde de işler ve çocuk
hastalığın nedeni (yanlış yaşam tarzı) olarak suçlanabilir --- bu,
ifşadan kaçınmanın içsel mantığını aydınlatır (Kim, 2022).
Düşük-kaynaklı bir bağlamda stigmanın önemli kaynaklarından biri
T1DM'nin bulaşıcı sanılması yanılgısıdır; akran stigması aidiyet
duygusunu aşındırırken, diyabetli aile üyelerinin ve akran/destek
gruplarının görünürlüğü kabullenme ve bağımsızlığı beslemektedir (Nsamba
ve ark., 2022). Genç erişkinlerde --- çocuk örnekleminden farklı bir
gelişimsel evre --- yürütülen bir nitel derleme, aynı gizleme/ifşa
ikileminin süreğenliğini ve sosyal ortamlarda insülin dozunun bilinçli
olarak atlanmasına dek varabildiğini bildirir; bu bulgu, çocukluk
döneminde gözlenen görünürlük kaygısının ileri evrelerdeki olası
uzantısına ilişkin bir bağlam sunar (Núñez-Baila ve ark., 2024).

\emph{Öne çıkan kodlar:} \texttt{KABULLENME\_NORMALLESME},
\texttt{RUTIN\_TAKIP}, \texttt{BESLENME\_KONTROL},
\texttt{OKUL\_SOSYAL\_UYUM}, \texttt{AKRAN\_CEVRE\_DESTEGI},
\texttt{COCUK\_OZERKLIK\_OZBAKIM}, \texttt{KARDES\_ILISKISI}.
\emph{Quote ID ankrajları:} \texttt{026\_t1dm\_child\_q001},
\texttt{026\_t1dm\_child\_q002}, \texttt{202\_t1dm\_child\_q001},
\texttt{202\_t1dm\_child\_q002}, \texttt{201\_t1dm\_child\_q001}.

\section{Tema 4 --- Aynı Evde Üç Farklı Deneyim}\label{tema-4}

\textbf{Ana sonuç.} Triadik okuma, T1DM'nin ailede tek bir ortak anlatı
üretmediğini; aynı olayın anne, hasta çocuk ve sağlıklı kardeş
tarafından farklı sorumluluk, risk ve adalet çerçeveleriyle
anlamlandırıldığını gösterir. Anne düzen/güvenlik dilini; hasta çocuk
beden/özerklik/sosyal görünürlük dilini; sağlıklı kardeş
kısıtlanma/gözetmenlik/adalet dilini öne çıkarır.

{\def\LTcaptype{none} % do not increment counter
\begin{longtable}[]{@{}
  >{\raggedright\arraybackslash}p{(\linewidth - 4\tabcolsep) * \real{0.3333}}
  >{\raggedright\arraybackslash}p{(\linewidth - 4\tabcolsep) * \real{0.3333}}
  >{\raggedright\arraybackslash}p{(\linewidth - 4\tabcolsep) * \real{0.3333}}@{}}
\toprule\noalign{}
\begin{minipage}[b]{\linewidth}\raggedright
Alt tema
\end{minipage} & \begin{minipage}[b]{\linewidth}\raggedright
Bilimsel netice
\end{minipage} & \begin{minipage}[b]{\linewidth}\raggedright
Tezdeki işlev
\end{minipage} \\
\midrule\noalign{}
\endhead
\bottomrule\noalign{}
\endlastfoot
Diyabetik düzenin gündelik hayatı yeniden kurması &
Yiyecek/zaman/planlama/ortak etkinlik düzeni değişir & T1DM'nin
aile-düzeyi yaşam organizasyonu \\
``Bir şey olacak'' korkusu ve sürekli tetikte olma & Risk beklentisi
rollere göre farklı yoğunluk/biçim & Kaygının rol-temelli dağılımı \\
İlgi, adalet ve kontrol gerilimi & Anne için koruma; kardeş için
görülmeme; çocuk için kontrol & H5 diadik tutarlılık için nitel köprü \\
Küçük bakıcılar ve hasta otoritesi & Kardeş görünmeyen yardımcı aktöre;
hasta çocuk hastalık statüsüyle farklı konuma & Kardeşlik dinamiklerinin
çok yönlü okuması \\
\end{longtable}
}

\textbf{Bilimsel yorum.} Tema 4, karma tezin niteliksel omurgasıdır.
Nicel H5 diadik tutarlılık bulguları sınırlı anne-çocuk örtüşmesi
gösterdiğinde, bu tema aynı bulguyu daha derin düzeyde açıklar:
tutarsızlık yalnız ölçüm hatası değil, aile içindeki rollerin aynı olayı
farklı ahlaki ve duygusal görevlerle taşıması olabilir.

Bu okuma, çok-bilgiverici (multi-informant) değerlendirme yazınının
merkezî bulgusuyla doğrudan örtüşür. Farklı türden bilgi vericiler arası
çapraz-uyum kuralsal olarak düşük-orta düzeydedir: klasik meta-analizde
benzer bilgi vericiler arası ortalama korelasyon r=.60 iken, farklı rol
taşıyan vericiler arası r=.28'e düşer --- bu örüntü bir ölçüm kusuru
değil, davranışın duruma-özgü (situational specificity) tezahürünün
beklenen yapısal sonucudur (Achenbach ve ark., 1987). Güncel kuramsal
çerçeve bu konumu ilerletir: bilgi verici farklılıkları sistematik
hataya indirgenemez, alan-ilişkili bilgi taşır ve klinik sonuçları
(tedavi yoğunluğu, uzunlamasına risk) öngörebilir; Operations Triad
Modeli uyumsuzluğu ``ıraksayan işlemler'' (diverging operations) olarak
konumlandırır (De Los Reyes ve ark., 2022; De Los Reyes, 2013). Bu
ilkenin anne-çocuk-kardeş triadına uygulanması --- kardeş, bu
kaynakların örneklemlerinde bir bilgi verici olarak yer almasa da ---
tezin kendi kuramsal çıkarımıdır.

Uyumsuzluğun ``aile içi rollerin farklı görevlerle taşınması''
biçimindeki aile-sistemleri okumamız ampirik dayanak da bulur. Pediatrik
astımda 279 çocuk-ebeveyn diyadında uyum zayıf-orta düzeyde (sınıf-içi
korelasyon katsayısı, ICC 0.32-0.47), çapraz-uyumsuzluk oranları ise
\%52.7-68.8 arasındadır; kritik olarak farklılığın hem büyüklüğü hem
yönü sosyodemografik/klinik değişkenlerden çok aile ilişki kalitesi ve
bakım yüküyle açıklanmıştır (Silva ve ark., 2015). T1DM bağlamında hem
ebeveyn hem çocuktaki psikolojik distres, diyabete-özgü aile
çatışmasının artışıyla yakından ilişkilidir (Williams ve ark., 2009).
Anne afektinin bu farklılığı biçimlendiren rolü, tezin maternal Beck
depresyon değişkeni açısından özellikle anlamlıdır: 73 anne-çocuk
diyadında maternal endişe ve depresyon, annenin çocuk anksiyetesini daha
düşük raporlamasını ve çapraz-uyumsuzluğun artmasını öngörmüştür
(Affrunti ve Woodruff-Borden, 2015). Son iki çalışma diyadik
tasarımlıdır (kardeşi içermez); triadik genelleme yorumsal bir
uzantıdır.

\emph{Öne çıkan kodlar:} \texttt{RUTIN\_TAKIP},
\texttt{KAYGI\_KIRILGANLIK}, \texttt{OFKE\_ADALETSIZLIK},
\texttt{KARDES\_GORUNMEZ\_YUK}, \texttt{KARDES\_ILISKISI},
\texttt{AILE\_ICI\_ADALET}, \texttt{COCUK\_OZERKLIK\_OZBAKIM}.

{\def\LTcaptype{none} % do not increment counter
\begin{longtable}[]{@{}
  >{\raggedright\arraybackslash}p{(\linewidth - 2\tabcolsep) * \real{0.5000}}
  >{\raggedright\arraybackslash}p{(\linewidth - 2\tabcolsep) * \real{0.5000}}@{}}
\toprule\noalign{}
\begin{minipage}[b]{\linewidth}\raggedright
Rol
\end{minipage} & \begin{minipage}[b]{\linewidth}\raggedright
Quote ID
\end{minipage} \\
\midrule\noalign{}
\endhead
\bottomrule\noalign{}
\endlastfoot
Anne & \texttt{011\_mother\_q006}, \texttt{020\_mother\_q009},
\texttt{026\_mother\_q008}, \texttt{026\_mother\_q009},
\texttt{011\_mother\_q007} \\
T1DM tanılı çocuk & \texttt{011\_t1dm\_child\_q004},
\texttt{020\_t1dm\_child\_q001}, \texttt{020\_t1dm\_child\_q002},
\texttt{026\_t1dm\_child\_q008}, \texttt{014\_t1dm\_child\_q002} \\
Sağlıklı kardeş & \texttt{011\_healthy\_sibling\_q001},
\texttt{020\_healthy\_sibling\_q001},
\texttt{019\_healthy\_sibling\_q002},
\texttt{019\_healthy\_sibling\_q003},
\texttt{201\_healthy\_sibling\_q001} \\
\end{longtable}
}

\chapter{Çapraz Bilimsel Neticeler}\label{sec-capraz}

\begin{enumerate}
\def\labelenumi{\arabic{enumi}.}
\tightlist
\item
  \textbf{T1DM aile-düzeyi bir düzenleme rejimi üretir.} Yemek, zaman,
  uyku, okul, etkinlik, sosyal görünürlük ve ebeveyn ilgisi yeniden
  düzenlenir; bu roller arası asimetri yaratır.
\item
  \textbf{Anne bakımı klinik ve ahlaki sorumluluk olarak içselleştirir.}
  ``İyi anne'' ölçütü ölçüm, izlem, kriz önleme ve adil davranma
  kapasitesiyle iç içe geçer.
\item
  \textbf{Sağlıklı kardeş yükü çoğu zaman sessizdir.} Geri
  çekilme/uyum/erteleme görünür uyum gibi okunabilir; görünmeyen yük
  taşıyabilir.
\item
  \textbf{T1DM tanılı çocuk normalleşme ve farklılık arasında müzakere
  eder.} İkilik aynı anlatıda bulunabilir.
\item
  \textbf{Koruma, kontrol ve adalet aynı davranışta birleşebilir.} İyi
  niyet deneyimsel sonuçla her zaman örtüşmez.
\item
  \textbf{Triadik farklılık bilimsel sonuçtur, hata değildir.} Amaç
  anlatıları tek mutabakata indirgemek değil, farklı öznellikleri
  birlikte okuyabilmektir.
\end{enumerate}

Bu çapraz neticeler aile-düzeyi bir kuramsal çerçeveye oturur.
Knafl-Deatrick Aile Yönetim Tarzı Çerçevesi (FMSF), kronik hastalıklı
çocuğa aile yanıtını durumun tanımı, yönetim davranışları ve algılanan
sonuçlar bileşenlerine ayırır (Knafl ve Deatrick, 2003; Knafl ve ark.,
2012); 575 ebeveyn ve 414 aileyi kapsayan bir çalışmada bu çerçeveden
türeyen dört yönetim deseni, aile-odaklı desenlerin daha iyi aile ve
çocuk işlevselliğiyle ilişkili olduğunu, tek ebeveynli ve düşük gelirli
ailelerin ise durum-odaklı (daha zorlu) desenlerde anlamlı olarak fazla
temsil edildiğini göstermiştir (Knafl ve ark., 2013). Ebeveyn tutumu ile
çocuğun klinik çıktısı arasındaki bağın bir örneği olarak, 251
ergen-ebeveyn diyadında çocuğun anksiyete/depresyon belirtileri ile
glisemik kontrol arasındaki ilişki tümüyle ebeveyn anksiyetesi üzerinden
aracılanmıştır (Silina ve ark., 2023); bu kesitsel bulgu nedensel bir
mekanizma iddiası değil, tetikte olma ikliminin aileye ve çocuğa
yayılımına ilişkin yorumsal bir zemindir.

\chapter{Negatif Vaka ve Sınırlayıcı Örüntüler}\label{sec-negatif}

Karşıt ve sınırlayıcı okumalar, dört makro temanın yorumuna gömülüdür
(ör. Tema 3'te normalleşmenin yük yokluğu anlamına gelmemesi; Tema 1'de
hastalığın kardeşler arası güç aracı olarak kullanımı). Tema 4 için
ayrıca sistematik bir negatif-vaka taraması yürütülmüş; anlatılardaki
karşıtlık ve sınırlayıcı okuma alanları belirlenmiştir.

\begin{tcolorbox}[enhanced jigsaw, arc=.35mm, bottomrule=.15mm, bottomtitle=1mm, breakable, colback=white, colbacktitle=quarto-callout-note-color!10!white, colframe=quarto-callout-note-color-frame, coltitle=black, left=2mm, leftrule=.75mm, opacityback=0, opacitybacktitle=0.6, rightrule=.15mm, title=\textcolor{quarto-callout-note-color}{\faInfo}\hspace{0.5em}{Not}, titlerule=0mm, toprule=.15mm, toptitle=1mm]

\textbf{Bilgi kutusu --- Negatif vaka analizi nedir?}

\textbf{Negatif vaka analizi}, ana temaya uymayan, onu sınırlayan veya
çelişen anlatıların bilinçli olarak aranıp raporlanmasıdır. Amaç,
örüntüyü ``bozan'' veriyi gizlemek yerine analize dahil ederek yorumun
sağlamlığını (credibility) sınamaktır; her makro temada en az bir
``farklı tezahür/sınırlayıcı örnek'' paragrafı COREQ Madde 31'i
güçlendirir (Tong ve ark., 2007). Aşağıdaki tablo bu koldaki karşıtlık
ve sınırlayıcı okuma alanlarını gösterir.

\end{tcolorbox}

{\def\LTcaptype{none} % do not increment counter
\begin{longtable}[]{@{}
  >{\raggedright\arraybackslash}p{(\linewidth - 2\tabcolsep) * \real{0.5000}}
  >{\raggedright\arraybackslash}p{(\linewidth - 2\tabcolsep) * \real{0.5000}}@{}}
\toprule\noalign{}
\begin{minipage}[b]{\linewidth}\raggedright
Alan
\end{minipage} & \begin{minipage}[b]{\linewidth}\raggedright
Bulgusal anlam
\end{minipage} \\
\midrule\noalign{}
\endhead
\bottomrule\noalign{}
\endlastfoot
Tema 4 içinde karşıtlık işaretleri & Tek yönlü tema anlatısı yerine
koşullu/sınırlı/çelişkili anlatılar bulunabilir \\
``Normal'' söylemi & Yük yokluğu olarak değil;
normalleştirme/dayanıklılık/performatif uyum olarak okunur \\
Aile 201 rol temsili & Bazı rol alanları sınırlı; düşük odak temsili
aktarılabilirlik tartışmasında raporlanır \\
\end{longtable}
}

Her makro temanın yorumu en az bir ``farklı tezahür/sınırlayıcı örnek''
içerir; her tema için ayrı sistematik bir negatif-vaka alt-bölümü bu
raporun kapsam sınırıdır (COREQ Madde 31; bkz.
Bölüm~\ref{sec-sinirliliklar}).

\chapter{Karma Tez Entegrasyonu}\label{sec-karma}

\begin{quote}
\textbf{Kanonik karma sentez pointer:} Doldurulmuş joint-display ve
meta-çıkarımın tek doğruluk kaynağı
\texttt{tez-yazim/05\_entegrasyon/karma-sentez-kanonik.md}'dir; buradaki
tablo yalnız nitel-kol bakışını sunar.
\end{quote}

Nitel bulgular nicel bulguların doğrulama aracı değildir; karma tezde üç
işlev görür: (i) ölçek sonuçlarını deneyimsel bağlamla açıklamak; (ii)
nicel fark yokluğu/zayıf etki durumlarında görünmeyen süreçleri
tartışmaya açmak; (iii) anne--hasta çocuk--kardeş algı ayrışmasının
neden bilimsel olarak anlamlı olduğunu göstermek. Yaklaşım,
eşzamanlı/convergent paralel karma yöntem tasarımının entegrasyon
disiplinini izler (Creswell ve Plano Clark, 2018); entegrasyon, joint
display ve convergence teknikleriyle yürütülür (Guetterman ve ark.,
2015; Fetters ve ark., 2013) (nicel kol köprüsü:
\texttt{t1dm-tez-rehberi/references/karma-yontem.md}).

{\def\LTcaptype{none} % do not increment counter
\begin{longtable}[]{@{}
  >{\raggedright\arraybackslash}p{(\linewidth - 4\tabcolsep) * \real{0.3333}}
  >{\raggedright\arraybackslash}p{(\linewidth - 4\tabcolsep) * \real{0.3333}}
  >{\raggedright\arraybackslash}p{(\linewidth - 4\tabcolsep) * \real{0.3333}}@{}}
\toprule\noalign{}
\begin{minipage}[b]{\linewidth}\raggedright
Nicel hat (CSR kanonik karar)
\end{minipage} & \begin{minipage}[b]{\linewidth}\raggedright
Nitel katkı
\end{minipage} & \begin{minipage}[b]{\linewidth}\raggedright
Karma kullanım dili
\end{minipage} \\
\midrule\noalign{}
\endhead
\bottomrule\noalign{}
\endlastfoot
H1 Çocuk algısı --- \textbf{doğrulanan pozitif}: DM çocukları EMBU-C
reddetme alt ölçeğinde daha yüksek (β ≈ 0,16 SD, BF₁₀ = 8,12);
alım-dönemi karışmasına karşı temkinli & Tema 3: hastalığı hem
normalleştirme hem yük olarak taşıma & Nitel bulgu farkı
\emph{doğrulamaz}; reddetme algısının çocuk tarafından nasıl
çerçevelendiğini verir (açıklayıcı genişleme) \\
H2 Kardeş ilişkisi --- \textbf{belirsiz / kanıt yetersizliği}: dört SRQ
alt ölçeğinde \textbar d\textbar{} \textless{} 0,20, FDR p
\textgreater{} 0,35 (\emph{aktif eşdeğerlik kanıtı değil}) & Tema 1 ve
4: yük genel puandan çok bakım rutini ve adalet algısına gömülü & Grup
farkı yokluğu kardeşin etkilenmediğini göstermez; nitel bulgu
sessiz/bağlamsal yük biçimlerini açıklar \\
H3 Anne öz-rapor --- \textbf{üç-katmanlı negatif kanıt}: kanıt
yetersizliği + TOST eşdeğerlik + Bayesçi (BF₁₀ = 0,17--0,25 H₀ lehine) &
Tema 2: tıbbi bakım rolünün normal ebeveynlik gibi içselleştirilmesi &
Ölçek eşdeğerliği psikolojik yük yokluğunu kanıtlamaz; içselleştirme
öz-raporu daraltabilir \\
H4 Beck → EMBU-P (SEM) --- \textbf{kısmen doğrulandı}:
sıcaklık/reddetme/karşılaştırma yolları anlamlı, aşırı koruma yolu değil
& Tema 2: suçluluk/tetikte olma/kaybetme korkusu ruhsal yükle ilişkili &
Kesitsel SEM; nedensel yol değil, uyum ve klinik yorum alanı \\
H5 Diadik tutarlılık --- \textbf{üçgenleme şartı karşılanmadı}: manifest
ICC dört alt ölçekte Kontrol \textgreater{} DM; latent DM \textgreater{}
Kontrol yalnız reddetmede/tek-strateji & Tema 4: aynı olayın farklı
anlamlarla taşınması & Algı ayrışması ölçüm hatası değil, rol-temelli
deneyim farkı \\
\end{longtable}
}

Kısaltmalar: EMBU-C / EMBU-P = algılanan ebeveynlik ölçeği çocuk /
ebeveyn formu; SRQ = Kardeş İlişkileri Anketi (KİA); BF₁₀ = Bayes
faktörü; SD = standart sapma; TOST = iki tek-yönlü eşdeğerlik testi; FDR
= yanlış keşif oranı; ICC = sınıf-içi korelasyon katsayısı.
\textbf{Nicel kararların kanonik kaynağı}
\texttt{docs/CLINICAL-STUDY-REPORT-FINAL} §2.3 yönetici özetidir;
doldurulmuş joint-display ve çapraz-kol meta-çıkarımların tek doğruluk
kaynağı \texttt{tez-yazim/05\_entegrasyon/karma-sentez-kanonik.md}'dir.

\textbf{Tartışma köprü cümleleri.} H2'de dört SRQ alt ölçeğinde grup
farkı için kanıt yetersizdir (aktif eşdeğerlik kanıtı değil); nitel
bulgular kardeş yükünün gündelik bakım rutinlerine, ilgi dağılımına ve
adalet algısına gömülü yaşanabileceğini gösterir; H3'te üç-katmanlı
negatif kanıt (Bayesçi + TOST eşdeğerlik dahil) anneliğin tıbbi bakıcı
rolüne kaymasının öz-raporda doğrudan belirmeyebileceğini açar; H5 Tema
4 ile birlikte yorumlandığında aynı aile olayının farklı
görev/risk/adalet çerçeveleriyle anlamlandırıldığını gösterir. Nitel
bulgular nedensel açıklama değil, deneyimsel bağlam ve yorumlayıcı zemin
olarak kullanılmalıdır. Aile yönetim tarzının nicel ölçüm karşılığı Aile
Yönetim Ölçütü (FaMM) ile sağlanabilir (Knafl ve ark., 2011); çocukların
yönetimi ``ben/onlar/biz'' ifadeleriyle anlattığı ve ebeveyn izlemini
kimi zaman güvensizlik olarak okuyabildiği gösterilmiştir (Beacham ve
Deatrick, 2019) --- bu da H5 kontrol-özerklik ekseninin nitel-nicel
yakınsaması için bir köprü sağlar.

\textbf{Ortak analiz için hazırlık.} İki kol eşzamanlı \textbf{paralel}
karma yöntem tasarımıyla yürütülür ve \emph{aynı bireyleri
örneklememiştir} (nitel: 7 aile / 21 görüşme; nicel: 241 aile / 482
katılımcı). Bu asimetri üç hazırlık ilkesini zorunlu kılar: (i)
\textbf{entegrasyon yorum düzeyindedir, istatistiksel düzeyde değil} ---
nitel bulgu nicel istatistiği doğrulamaz, bağlamsal anlamını
derinleştirir; bu nedenle ``birbirini doğrular'' yerine ``tamamlar /
açıklayıcı bağlam sağlar'' retoriği kullanılır (Fetters ve ark., 2013).
(ii) \textbf{Aktarılabilirlik analitik genellemeyle kurulur},
istatistiksel genellemeyle değil; 7 ailelik nitel bağlamın 241 ailelik
örnekleme taşınması analitik çıkarım yoludur. (iii) Her H-hattının nicel
ankrajı (\texttt{\#h1-karar\ …\ \#h5-karar}) ile nitel tema ankrajı
(\texttt{\#tema-1\ …\ \#tema-4})
\texttt{tez-yazim/05\_entegrasyon/karma-kanit-ledgeri.tsv} üzerinden
eşlenir; doldurulmuş joint-display ve dört çapraz-kol meta-çıkarım
\texttt{karma-sentez-kanonik.md}'de tutulur. Bu paralel-örneklem sınırı
Sınırlılıklar (bkz. Bölüm~\ref{sec-sinirliliklar}) ve karma-özel
sınırlılık kesitinde şeffaf raporlanır.

\chapter{COREQ, Trustworthiness ve Değerlendirme}\label{sec-raporlama}

\section{COREQ Durumu}\label{coreq-durumu}

\begin{tcolorbox}[enhanced jigsaw, arc=.35mm, bottomrule=.15mm, bottomtitle=1mm, breakable, colback=white, colbacktitle=quarto-callout-note-color!10!white, colframe=quarto-callout-note-color-frame, coltitle=black, left=2mm, leftrule=.75mm, opacityback=0, opacitybacktitle=0.6, rightrule=.15mm, title=\textcolor{quarto-callout-note-color}{\faInfo}\hspace{0.5em}{Not}, titlerule=0mm, toprule=.15mm, toptitle=1mm]

\textbf{Bilgi kutusu --- COREQ nedir ve ``0 eksik'' neyi ifade eder?}

COREQ, nitel görüşme/odak grup çalışmalarının \textbf{raporlanmasını}
standartlaştıran 32 maddelik bir kontrol listesidir; üç alanı kapsar:
araştırma ekibi ve refleksivite, çalışma tasarımı, analiz ve bulgular
(Tong ve ark., 2007). COREQ bir \textbf{raporlama} ölçütüdür, bir
\emph{yöntem zorunluluğu} değil: bir madde, ilgili uygulamanın yapılıp
yapılmadığının \textbf{şeffaf biçimde raporlanmasıyla} karşılanır. Bu
nedenle transkript iadesi (Madde 22) ve resmî katılımcı doğrulaması
(Madde 27) bu çalışmada \emph{uygulanmamış} olsa da, uygulanmama kararı
ve yerine konan görüşme-içi özetleme gerekçesiyle açıkça
raporlandığından COREQ mantığında ``karşılanmış'' sayılır. Yani ``0
eksik'', tüm uygulamaların \emph{yapıldığı} değil, tüm maddelerin
\emph{şeffaf raporlandığı} anlamına gelir.

\end{tcolorbox}

{\def\LTcaptype{none} % do not increment counter
\begin{longtable}[]{@{}lrrrr@{}}
\toprule\noalign{}
Domain & Tam & Kısmi & Eksik & Toplam \\
\midrule\noalign{}
\endhead
\bottomrule\noalign{}
\endlastfoot
Araştırma ekibi ve refleksivite & 7 & 0 & 0 & 7 \\
Çalışma tasarımı & 14 & 1 & 0 & 15 \\
Analiz ve bulgular & 9 & 1 & 0 & 10 \\
\textbf{Toplam} & \textbf{30} & \textbf{2} & \textbf{0} & \textbf{32} \\
\end{longtable}
}

Kısmi maddeler: \textbf{Madde 16} (ayrı formal pilot görüşme
yürütülmemiştir; rehber alanyazın ve klinik deneyimle geliştirilip
sahada iyileştirilmiştir, şeffaf raporlanmıştır) ve \textbf{Madde 31}
(karşıt/ sınırlayıcı okumalar tema yorumlarına gömülüdür; her tema için
ayrı sistematik alt-bölüm kapsam sınırıdır) (Tong ve ark., 2007).
Raporlama ayrıca SRQR (O'Brien ve ark., 2014) ve JARS-Qual (Levitt ve
ark., 2018) ölçütleriyle çapraz denetlenir. Transkript iadesi (Madde 22)
ve resmî katılımcı doğrulaması (member checking / respondent validation;
Madde 27) (Birt ve ark., 2016), yukarıdaki bilgi kutusunda açıklandığı
üzere uygulanmayıp gerekçesiyle raporlandıkları için ``karşılanmış''
sayılmıştır; bu maddelerin uygulama sınırı Sınırlılıklar bölümünde
ayrıca ele alınmıştır (bkz. Bölüm~\ref{sec-sinirliliklar}). Kaynak
metodoloji dosyası (\texttt{coreq\_32\_completed.md}) madde-madde
tablosu ve domain özeti bu tabloyla \textbf{tek dildedir}: 30 tam / 2
kısmi / 0 eksik, kısmiler yalnız Madde 16 (pilot) ve Madde 31 (negatif
vaka).

\section{Trustworthiness (Lincoln-Guba) ve Değerlendirme
Ölçütleri}\label{trustworthiness-lincoln-guba-ve-deux11ferlendirme-uxf6luxe7uxfctleri}

Niteliksel kalite, pozitivist iç/dış geçerlik yerine Lincoln-Guba
\textbf{trustworthiness} paradigmasıyla kurulur (Lincoln ve Guba, 1985);
kalite değerlendirmesi Tracy'nin ``big-tent'' ölçütleriyle desteklenir
(Tracy, 2010).

\begin{tcolorbox}[enhanced jigsaw, arc=.35mm, bottomrule=.15mm, bottomtitle=1mm, breakable, colback=white, colbacktitle=quarto-callout-note-color!10!white, colframe=quarto-callout-note-color-frame, coltitle=black, left=2mm, leftrule=.75mm, opacityback=0, opacitybacktitle=0.6, rightrule=.15mm, title=\textcolor{quarto-callout-note-color}{\faInfo}\hspace{0.5em}{Not}, titlerule=0mm, toprule=.15mm, toptitle=1mm]

\textbf{Bilgi kutusu --- Trustworthiness: nitel araştırmada
``geçerlik''in karşılığı}

Lincoln ve Guba, pozitivist iç/dış geçerlik yerine nitel araştırmaya
özgü dört \textbf{güvenilirlik (trustworthiness)} ölçütü önerir (Lincoln
ve Guba, 1985):

\begin{itemize}
\tightlist
\item
  \textbf{Credibility (inandırıcılık)} ↔ iç geçerlik: bulgular veriyi
  sadık biçimde temsil ediyor mu?
\item
  \textbf{Transferability (aktarılabilirlik)} ↔ dış geçerlik: bulgular
  hangi bağlamlara taşınabilir?
\item
  \textbf{Dependability (tutarlılık)} ↔ güvenirlik: süreç izlenebilir ve
  tutarlı mı?
\item
  \textbf{Confirmability (doğrulanabilirlik)} ↔ nesnellik: bulgular
  araştırmacı önyargısından çok veriden mi doğuyor?
\end{itemize}

Aşağıdaki tablo her ölçütün bu çalışmadaki operasyonel dayanağını verir.

\end{tcolorbox}

{\def\LTcaptype{none} % do not increment counter
\begin{longtable}[]{@{}
  >{\raggedright\arraybackslash}p{(\linewidth - 2\tabcolsep) * \real{0.5000}}
  >{\raggedright\arraybackslash}p{(\linewidth - 2\tabcolsep) * \real{0.5000}}@{}}
\toprule\noalign{}
\begin{minipage}[b]{\linewidth}\raggedright
Ölçüt
\end{minipage} & \begin{minipage}[b]{\linewidth}\raggedright
Bu çalışmadaki dayanak
\end{minipage} \\
\midrule\noalign{}
\endhead
\bottomrule\noalign{}
\endlastfoot
Credibility & Triadik çoklu bilgi verici yapı, \texttt{quote\_id}
ankrajlı yoğun kanıt, görüşme-içi özetleme, critical friend
tartışmaları \\
Dependability & Audit trail, kanonik codebook (codebook\_v2),
metodolojik karar kayıtları \\
Confirmability & Reflexive memo, positionality (OM/BA), quote integrity
raporu, AI-use ledger \\
Transferability & Bilgi gücü gerekçesi, bağlamın yoğun betimlemesi
(thick description (Ponterotto, 2006)), sınırlılıkların (ve Aile 201)
açık raporlanması \\
\end{longtable}
}

\begin{tcolorbox}[enhanced jigsaw, arc=.35mm, bottomrule=.15mm, bottomtitle=1mm, breakable, colback=white, colbacktitle=quarto-callout-note-color!10!white, colframe=quarto-callout-note-color-frame, coltitle=black, left=2mm, leftrule=.75mm, opacityback=0, opacitybacktitle=0.6, rightrule=.15mm, title=\textcolor{quarto-callout-note-color}{\faInfo}\hspace{0.5em}{Not}, titlerule=0mm, toprule=.15mm, toptitle=1mm]

\textbf{Bilgi kutusu --- RTA'da neden ikinci-kodlayıcı uzlaşımı (κ / α)
hesaplanmaz?}

Kodlama-güvenirliği (coding reliability) yaklaşımları, iki kodlayıcının
aynı kodu üretmesini bir nesnellik kanıtı sayar ve Cohen's κ /
Krippendorff α / Gwet AC1 gibi uzlaşım katsayıları hesaplar. Refleksif
TA'nın epistemolojisinde ise kodlama \textbf{yorumsal bir kaynaktır}:
amaç kodlayıcıların örtüşmesi değil, veriye derinlemesine nüfuz
etmektir. Bu nedenle uzlaşım katsayısı hesaplamak paradigmatik olarak
\emph{uygun değildir}; güvenilirlik bunun yerine \textbf{critical
friend} tartışmaları, refleksivite ve audit trail ile kurulur (Braun ve
Clarke, 2019).

\end{tcolorbox}

\textbf{Inter-coder reliability neden yok.} RTA epistemolojisinde
kodlama yorumsal bir kaynaktır; iki kodlayıcının aynı kodu üretmesi amaç
değildir. Bu nedenle Cohen's κ / Krippendorff α / Gwet AC1
\textbf{hesaplanmamıştır} (Braun ve Clarke, 2019); güvenilirlik critical
friend, refleksivite ve audit trail üzerinden kurulur. (Nicel skill'in
karma-yöntem dosyasındaki coding-reliability/AC1 önerisi bu kanonik RTA
kola uygulanmaz.)

\section{Alıntı
Bütünlüğü}\label{alux131ntux131-buxfctuxfcnluxfcux11fuxfc}

Bu belge doğrudan alıntı metni taşımadığı için alıntı kanıtı yalnız
\texttt{quote\_id} düzeyindedir. Mevcut \texttt{quotes\_used.csv}
ID-only yapıdadır; \texttt{dmnitel\ check-quotes} aracının beklediği
\texttt{quote\_text\_used} alanı bulunmadığından, doğrudan alıntı metni
bütünlüğü bu dosya üzerinden kapanmış kabul edilmez. Final tez veya
makale metnine verbatim alıntı eklendiğinde, her alıntı temizlenmiş tez
metniyle bire bir yeniden denetlenmelidir. Alıntının katılımcının kendi
sözlerine sadakati ile anonimleştirme beklentisi, niteliksel raporlama
yazınında yerleşik bir ölçüttür (Corden ve Sainsbury, 2006; Levitt ve
ark., 2018; Tong ve ark., 2007). İzin verilen müdahaleler: kimlik
maskeleme (\texttt{{[}isim{]}}, \texttt{{[}okul\ adı{]}},
\texttt{{[}şehir{]}}), kısaltma (\texttt{{[}...{]}}), köşeli parantezde
sınırlı bağlam ve okunabilirlik için asgari düzenleme (dolgu
sözcük/tekrar ayıklama). Yasak: \textbf{anlamı değiştiren}
akademikleştirme/yeniden yazma, özet metni alıntı gibi sunma, alıntıdan
tema üretimini kesin hüküm gibi sunma.

\chapter{Sınırlılıklar}\label{sec-sinirliliklar}

\begin{enumerate}
\def\labelenumi{\arabic{enumi}.}
\tightlist
\item
  Örneklem tek merkezli ve küçük bir triadik örneklemdir; RTA + bilgi
  gücü açısından kabul edilebilir olsa da istatistiksel genelleme
  sağlamaz.
\item
  Aileler arası veri zenginliği eşit değildir; Aile 201 düşük odak
  temsilindedir.
\item
  Formal post-hoc member checking yapılmamış; görüşme-içi özetleme +
  katılımcı onayı kullanılmıştır.
\item
  Ayrı bir formal pilot görüşme yürütülmemiştir; rehber alanyazın ve
  klinik deneyimle geliştirilip sahada iyileştirilmiştir (COREQ Madde
  16).
\item
  Babalar/diğer bakım verenler triadın içinde değildir; aile sistemi
  anne--hasta--kardeş ekseninden okunmuştur.
\item
  BA'nın gözlemci olarak odada bulunması katılımcı anlatısını
  şekillendirmiş olabilir.
\item
  OM ve BA'nın klinik disiplin ortaklığı bazı klinik-dışı yapısal
  okumaları sınırlamış olabilir.
\item
  Codebook'un beş kavramsal kodu seçilmiş illüstratif alıntı havuzunda
  temsil edilmez; illüstratif havuz tüketici değil, purposive bir
  seçkidir.
\item
  Kodlama araştırmacı-yorumlu ve özyinelidir; RTA epistemolojisi gereği
  pozitivist anlamda kilitli bir kod kümesi olarak sunulmaz.
\item
  Sistematik negatif-vaka taraması Tema 4 için yürütülmüştür; karşıt
  okumalar diğer temaların yorumuna gömülüdür, ancak her tema için ayrı
  sistematik alt-bölüm bu raporun kapsamı dışındadır (COREQ Madde 31).
\end{enumerate}

\chapter{AI-Reliability ve Yönetişim Katmanı}\label{sec-yonetisim}

Niteliksel kol, üretilen çıktının güvenilirliğini artırmak ve
denetlenebilir kılmak için çok katmanlı bir AI-reliability yönetişimi
taşır. Bu katman bilimsel bulgu üretmez; \textbf{bulgu ve metnin
bütünlüğünü denetler}.

\begin{itemize}
\tightlist
\item
  \textbf{Yerel toolkit (\texttt{dmnitel}):} codebook lint, triadik
  matris, COREQ audit, quote integrity, negatif vaka bulucu, AI-use
  ledger --- harici API'siz, ham veri connector'a gitmez.
\item
  \textbf{Reliability hattı:} claim-grounding orkestrasyonu, Lynx
  halüsinasyon eval, RAGAS faithfulness, semantic entropy tutarlılık
  sinyali; promptfoo eval + red-team.
\item
  \textbf{Yönetişim:} NIST--ISO/IEC 42001 eşlemesi;
  \texttt{t1dm-qual-ai-audit} referanslı-bölüm denetim kapısı.
\item
  \textbf{AI-use ledger:} her LLM/MCP kullanımı tarih, model, amaç, veri
  türü, ham-veri/kimlikleyici bayrağı ve araştırmacı kararı ile
  kaydedilir.
\item
  \textbf{Manüskript denetimi:} bu belge kapanışta \texttt{sci-audit}
  yedi ekseninden (A referans, B claim grounding, C istatistik, D
  halüsinasyon, \textbf{E COREQ/JARS-Qual}, F AI-şeffaflık, G Türkçe
  imla) ve bağımsız Galileo three-tier (GPT-5.4 ikinci görüş) pass'inden
  geçirilir.
\end{itemize}

\chapter{Kaynak Dosya Haritası}\label{sec-kaynak-harita}

{\def\LTcaptype{none} % do not increment counter
\begin{longtable}[]{@{}
  >{\raggedright\arraybackslash}p{(\linewidth - 2\tabcolsep) * \real{0.5000}}
  >{\raggedright\arraybackslash}p{(\linewidth - 2\tabcolsep) * \real{0.5000}}@{}}
\toprule\noalign{}
\begin{minipage}[b]{\linewidth}\raggedright
Kaynak
\end{minipage} & \begin{minipage}[b]{\linewidth}\raggedright
İşlev
\end{minipage} \\
\midrule\noalign{}
\endhead
\bottomrule\noalign{}
\endlastfoot
\texttt{03\_analysis/codebook/codebook\_v2.md} & 6 journal / 4 tez tema
mapping (kanonik) \\
\texttt{03\_analysis/codebook/codebook\_v3.csv} & Kanonik codebook\_v2
tema tabanının denetlenebilir CSV temsili \\
\texttt{06\_manuscript\_outputs/quotes\_used.csv} & \texttt{quote\_id}
kayıt dosyası; doğrudan alıntı metni içermez (kanıt \texttt{quote\_id}
düzeyinde ankrajlanır) \\
\texttt{03\_analysis/methodology/A1\_information\_power.md} & Bilgi gücü
/ doygunluk reframe dayanağı \\
\texttt{03\_analysis/methodology/A9\_triadic\_methodology\_literature.md}
& Triadik/multi-informant atıf dayanağı \\
\texttt{03\_analysis/methodology/coreq\_32\_completed.md} & COREQ 32
madde durumu \\
\texttt{03\_analysis/methodology/audit\_trail.md} & Metodolojik karar
günlüğü \\
\texttt{03\_analysis/methodology/positionality\_OM.md},
\texttt{positionality\_BA.md} & Refleksivite beyanları \\
\texttt{03\_analysis/methodology/llm\_use\_statement.md} & AI/LLM
kullanım beyanı \\
\texttt{03\_analysis/reflexive/journal\_excerpts.md} & Refleksif
günlük \\
\texttt{07\_reports/quote\_integrity\_report.md} & Alıntı bütünlüğü
denetim kaydı \\
\texttt{07\_reports/negative\_case\_report\_tema\_4\_triadic.md} & Tema
4 sistematik negatif vaka analizi \\
\texttt{qualitative\_canonical\_results\_report.md} & Bu QMD'nin mekanik
Markdown kopyası; tek başına ayrı kanonik otorite değildir \\
\end{longtable}
}

\chapter{Kalite Güvencesi ve Provenans}\label{sec-dogrulama}

Bu belgenin kalite güvencesi ve kanıt temeli şu katmanlara dayanır:

\begin{itemize}
\tightlist
\item
  \textbf{Kanıt temeli.} Alıntı kanıtı \texttt{quote\_id} düzeyinde
  ankrajlanır; verbatim alıntılar tez/makale metnine temizlenmiş kaynak
  metinden \texttt{quote\_id} eşlemesiyle bire bir alınır (KVKK + alıntı
  bütünlüğü ölçütü). Aile, katılımcı, tema, kod ve COREQ sayıları
  kanonik codebook ve methodology kaynaklarıyla tutarlıdır
  (\texttt{codebook\_v2.md}, \texttt{coreq\_32\_completed.md},
  \texttt{A1\_information\_power.md}).
\item
  \textbf{Referans tek-kaynağı.} Tüm \texttt{{[}@key{]}} atıfları tek
  \texttt{references/references.bib} (Zotero koleksiyonu
  \textbf{9ZFDHMZA ``T1DM Thesis''} ile mutabık) üzerinden çözülür; bib
  hijyeni \texttt{scripts/util/bib\_hygiene.py\ all} ile HARD=0
  doğrulanır.
\item
  \textbf{Dış-kanıt.} Metodoloji atıfları evidentia bağlı
  connector'larıyla (OpenAlex + PubMed/EPMC) çözülüp doğrulanmıştır;
  bulunamayan kaynak \texttt{gap} olarak işlenir, uydurma atıf
  üretilmez.
\item
  \textbf{Metin denetimi.} Belge, \texttt{sci-audit} yedi ekseni (A
  referans, B claim grounding, C istatistik, D halüsinasyon, E
  COREQ/JARS-Qual, F AI-şeffaflık, G Türkçe imla) ve bağımsız Galileo
  three-tier judge denetiminden geçirilir.
\end{itemize}

\chapter{Nihai Bilimsel Kullanım Cümlesi}\label{sec-nihai}

Bu niteliksel kol, T1DM'nin aile içinde yalnız hastalık yönetimi değil;
rol, adalet, bakım emeği, görünürlük ve özerklik düzenini yeniden kuran
bir aile deneyimi olduğunu göstermektedir. Anne, hasta çocuk ve sağlıklı
kardeş aynı klinik durumu farklı biçimlerde taşımakta; bu farklılıklar
özellikle kardeşin görünmeyen yükü, annenin tıbbi bakıcı rolüne kayması,
çocuğun normalleşme--farklılık müzakeresi ve triadik adalet/kontrol
gerilimleri üzerinden görünür olmaktadır. Karma tezde bu sonuçlar, nicel
bulguları ikame eden veya doğrulayan kanıtlar olarak değil; aile içi
algı farklılıklarını ve ölçekte görünmeyebilecek deneyim süreçlerini
açıklayan niteliksel bağlam olarak kullanılmalıdır.

\chapter*{Kaynakça}\label{kaynakuxe7a}
\addcontentsline{toc}{chapter}{Kaynakça}

\protect\phantomsection\label{refs}
\begin{CSLReferences}{1}{0}
\bibitem[\citeproctext]{ref-achenbach1987crossinformant}
Achenbach TM, McConaughy SH, Howell CT. Child/adolescent behavioral and
emotional problems: Implications of cross-informant correlations for
situational specificity. Psychological Bulletin. 1987;101(2):213-232.
\url{https://doi.org/10.1037/0033-2909.101.2.213}

\bibitem[\citeproctext]{ref-adams1991siblings}
Adams R, Peveler RC, Stein A, Dunger DB. Siblings of Children with
Diabetes: Involvement, Understanding and Adaptation. Diabetic Medicine.
1991;8(9):855-859.
\url{https://doi.org/10.1111/j.1464-5491.1991.tb02124.x}

\bibitem[\citeproctext]{ref-affrunti2015maternal}
Affrunti NW, Woodruff-Borden J. The effect of maternal psychopathology
on parent--child agreement of child anxiety symptoms: A hierarchical
linear modeling approach. Journal of Anxiety Disorders. 2015;32:56-65.
\url{https://doi.org/10.1016/j.janxdis.2015.03.010}

\bibitem[\citeproctext]{ref-azimi2024caregiver}
Azimi T, Johnson J, Campbell SM, Montesanti S. Caregiver burden among
parents of children with type 1 diabetes: A qualitative scoping review.
Heliyon. 2024;10(6):e27539.
\url{https://doi.org/10.1016/j.heliyon.2024.e27539}

\bibitem[\citeproctext]{ref-barnard2010fear}
Barnard K, Thomas S, Royle P, Noyes K, Waugh N. Fear of hypoglycaemia in
parents of young children with type 1 diabetes: a systematic review. BMC
Pediatrics. 2010;10(1):50. \url{https://doi.org/10.1186/1471-2431-10-50}

\bibitem[\citeproctext]{ref-beacham2019children}
Beacham BL, Deatrick JA. Adapting the Family Management Styles Framework
to Include Children. Journal of Pediatric Nursing. 2019;45:26-36.
\url{https://doi.org/10.1016/j.pedn.2018.12.006}

\bibitem[\citeproctext]{ref-birt2016member}
Birt L, Scott S, Cavers D, Campbell C, Walter F. Member Checking: A Tool
to Enhance Trustworthiness or Merely a Nod to Validation? Qualitative
Health Research. 2016;26(13):1802-1811.
\url{https://doi.org/10.1177/1049732316654870}

\bibitem[\citeproctext]{ref-braunClarke2006thematic}
Braun V, Clarke V. Using thematic analysis in psychology. Qualitative
Research in Psychology. 2006;3(2):77-101.
\url{https://doi.org/10.1191/1478088706qp063oa}

\bibitem[\citeproctext]{ref-braunClarke2019reflexive}
Braun V, Clarke V. Reflecting on reflexive thematic analysis.
Qualitative Research in Sport, Exercise and Health. 2019;11(4):589-597.
\url{https://doi.org/10.1080/2159676X.2019.1628806}

\bibitem[\citeproctext]{ref-braunClarke2019saturate}
Braun V, Clarke V. To saturate or not to saturate? Questioning data
saturation as a useful concept for thematic analysis and sample-size
rationales. Qualitative Research in Sport, Exercise and Health.
2021;13(2):201-216. \url{https://doi.org/10.1080/2159676X.2019.1704846}

\bibitem[\citeproctext]{ref-cao2021family}
Cao VT, Anderson BJ, Eshtehardi SS, McKinney BM, Thompson DI, Marrero DG
ve ark. "We are a family with diabetes": Parent perspectives on siblings
of youth with type 1 diabetes. Families, Systems, \& Health.
2021;39(2):306-315. \url{https://doi.org/10.1037/fsh0000612}

\bibitem[\citeproctext]{ref-carlsund2025stress}
Carlsund Å, Olsson S, Hörnsten Å. Stress and Burden Experienced by
Parents of Children with Type 1 Diabetes: A Qualitative Content Analysis
Interview Study. Children. 2025;12(8):984.
\url{https://doi.org/10.3390/children12080984}

\bibitem[\citeproctext]{ref-corden2006quotations}
Corden A, Sainsbury R. Exploring {``Quality''}: Research Participants'
Perspectives on Verbatim Quotations. International Journal of Social
Research Methodology. 2006;9(2):97-110.
\url{https://doi.org/10.1080/13645570600595264}

\bibitem[\citeproctext]{ref-creswellPlanoClark2018}
Creswell JW, Plano Clark VL. Designing and Conducting Mixed Methods
Research. 2018.

\bibitem[\citeproctext]{ref-delosReyes2013strategic}
De Los Reyes A. Strategic objectives for improving understanding of
informant discrepancies in developmental psychopathology research.
Development and Psychopathology. 2013;25(3):669-682.
\url{https://doi.org/10.1017/S0954579413000096}

\bibitem[\citeproctext]{ref-deLosReyes2015}
De Los Reyes A, Augenstein TM, Wang M, Thomas SA, Drabick DAG, Burgers
DE ve ark. The Validity of the Multi-Informant Approach to Assessing
Child and Adolescent Mental Health. Psychological Bulletin.
2015;141(4):858-900. \url{https://doi.org/10.1037/a0038498}

\bibitem[\citeproctext]{ref-delosReyes2022discrepancies}
De Los Reyes A, Tyrell FA, Watts AL, Asmundson GJG. Conceptual,
methodological, and measurement factors that disqualify use of
measurement invariance techniques to detect informant discrepancies in
youth mental health assessments. Frontiers in Psychology.
2022;13:931296. \url{https://doi.org/10.3389/fpsyg.2022.931296}

\bibitem[\citeproctext]{ref-deatrick1999normalization}
Deatrick JA, Knafl KA, Murphy-Moore C. Clarifying the Concept of
Normalization. Image: The Journal of Nursing Scholarship.
1999;31(3):209-214.
\url{https://doi.org/10.1111/j.1547-5069.1999.tb00482.x}

\bibitem[\citeproctext]{ref-elhabashy2023siblings}
Elhabashy SA, Elhenawy YI, Hassan HA, Abdelmageed RI. Siblings of
children with type 1 diabetes mellitus: Psychosocial health, coping.
Journal of Pediatric Nursing. 2023;73:e516-e524.
\url{https://doi.org/10.1016/j.pedn.2023.10.024}

\bibitem[\citeproctext]{ref-fetters2013integration}
Fetters MD, Curry LA, Creswell JW. Achieving Integration in Mixed
Methods Designs---Principles and Practices. Health Services Research.
2013;48(6pt2):2134-2156. \url{https://doi.org/10.1111/1475-6773.12117}

\bibitem[\citeproctext]{ref-guetterman2015jointDisplay}
Guetterman TC, Fetters MD, Creswell JW. Integrating Quantitative and
Qualitative Results in Health Science Mixed Methods Research Through
Joint Displays. Annals of Family Medicine. 2015;13(6):554-561.
\url{https://doi.org/10.1370/afm.1865}

\bibitem[\citeproctext]{ref-haghighiMoghadam2022mothers}
Haghighi Moghadam Y, Zeinaly Z, Alhani F. How mothers of a child with
type 1 diabetes cope with the burden of care: a qualitative study. BMC
Endocrine Disorders. 2022;22(1):129.
\url{https://doi.org/10.1186/s12902-022-01045-z}

\bibitem[\citeproctext]{ref-kelada2022siblings}
Kelada L, Wakefield CE, Drew D, Ooi CY, Palmer EE, Bye A ve ark.
Siblings of young people with chronic illness: Caring responsibilities
and psychosocial functioning. Journal of Child Health Care.
2022;26(4):581-596. \url{https://doi.org/10.1177/13674935211033466}

\bibitem[\citeproctext]{ref-kim2022illness}
Kim JE. Illness Experiences of Adolescents with Type 1 Diabetes. Journal
of Diabetes Research. 2022;2022:3117253.
\url{https://doi.org/10.1155/2022/3117253}

\bibitem[\citeproctext]{ref-knafl2003fmsf}
Knafl KA, Deatrick JA. Further Refinement of the Family Management Style
Framework. Journal of Family Nursing. 2003;9(3):232-256.
\url{https://doi.org/10.1177/1074840703255435}

\bibitem[\citeproctext]{ref-knafl2012continued}
Knafl KA, Deatrick JA, Havill NL. Continued Development of the Family
Management Style Framework. Journal of Family Nursing. 2012;18(1):11-34.
\url{https://doi.org/10.1177/1074840711427294}

\bibitem[\citeproctext]{ref-knafl2013patterns}
Knafl KA, Deatrick JA, Knafl GJ, Gallo AM, Grey M, Dixon J. Patterns of
Family Management of Childhood Chronic Conditions and Their Relationship
to Child and Family Functioning. Journal of Pediatric Nursing.
2013;28(6):523-535. \url{https://doi.org/10.1016/j.pedn.2013.03.006}

\bibitem[\citeproctext]{ref-knafl2011famm}
Knafl K, Deatrick JA, Gallo A, Dixon J, Grey M, Knafl G ve ark.
Assessment of the Psychometric Properties of the Family Management
Measure. Journal of Pediatric Psychology. 2011;36(5):494-505.
\url{https://doi.org/10.1093/jpepsy/jsp034}

\bibitem[\citeproctext]{ref-kobos2023loneliness}
Kobos E, Rojkowska S, Szewczyk A, Dziedzic B. Burden of care and a sense
of loneliness in caregivers of children with type 1 diabetes: a
cross-sectional study. BioPsychoSocial Medicine. 2023;17(1):34.
\url{https://doi.org/10.1186/s13030-023-00291-4}

\bibitem[\citeproctext]{ref-levitt2018jarsQual}
Levitt HM, Bamberg M, Creswell JW, Frost DM, Josselson R, Suárez-Orozco
C. Journal article reporting standards for qualitative primary,
qualitative meta-analytic, and mixed methods research in psychology:
{The} {APA} {Publications} and {Communications} {Board} task force
report. American Psychologist. 2018;73(1):26-46.
\url{https://doi.org/10.1037/amp0000151}

\bibitem[\citeproctext]{ref-lincolnGuba1985}
Lincoln YS, Guba EG. Naturalistic Inquiry. 1985.

\bibitem[\citeproctext]{ref-malterud2016informationPower}
Malterud K, Siersma VD, Guassora AD. Sample Size in Qualitative
Interview Studies: Guided by Information Power. Qualitative Health
Research. 2016;26(13):1753-1760.
\url{https://doi.org/10.1177/1049732315617444}

\bibitem[\citeproctext]{ref-uganda2022lived}
Nsamba J, Nabirye G, Hense S, Drenos F, Mathews E. Lived Experiences of
Newly Diagnosed Type 1 Diabetes Mellitus Children and Adolescents in
Uganda. Journal of Multidisciplinary Healthcare. 2022;15:2647-2665.
\url{https://doi.org/10.2147/JMDH.S389265}

\bibitem[\citeproctext]{ref-emergingadults2024lifestyle}
Núñez-Baila M-Á, Gómez-Aragón A, Marques-Silva A-M, González-López JR.
Lifestyle in Emerging Adults with Type 1 Diabetes Mellitus: A
Qualitative Systematic Review. Healthcare. 2024;12(3):309.
\url{https://doi.org/10.3390/healthcare12030309}

\bibitem[\citeproctext]{ref-obrien2014srqr}
O'Brien BC, Harris IB, Beckman TJ, Reed DA, Cook DA. Standards for
{Reporting} {Qualitative} {Research}. Academic Medicine.
2014;89(9):1245-1251. \url{https://doi.org/10.1097/acm.0000000000000388}

\bibitem[\citeproctext]{ref-palmer2022kenya}
Palmer T, Waliaula C, Shannon G, Salustri F, Grewal G, Chelagat W ve
ark. Understanding the Lived Experience of Children With Type 1 Diabetes
in Kenya: Daily Routines and Adaptation Over Time. Qualitative Health
Research. 2022;32(1):145-158.
\url{https://doi.org/10.1177/10497323211049775}

\bibitem[\citeproctext]{ref-ponterotto2006thick}
Ponterotto JG. Brief Note on the Origins, Evolution, and Meaning of the
Qualitative Research Concept Thick Description. The Qualitative Report.
2006;11(3):538-549. \url{https://doi.org/10.46743/2160-3715/2006.1666}

\bibitem[\citeproctext]{ref-rankin2014pathways}
Rankin D, Harden J, Waugh N, Noyes K, Barnard KD, Stephen J ve ark.
Pathways to diagnosis: a qualitative study of the experiences and
emotional reactions of parents of children diagnosed with type 1
diabetes. Pediatric Diabetes. 2014;15(8):591-598.
\url{https://doi.org/10.1111/pedi.12124}

\bibitem[\citeproctext]{ref-robinson1993normalization}
Robinson CA. Managing Life with a Chronic Condition: The Story of
Normalization. Qualitative Health Research. 1993;3(1):6-28.
\url{https://doi.org/10.1177/104973239300300102}

\bibitem[\citeproctext]{ref-saunders2018saturation}
Saunders B, Sim J, Kingstone T, Baker S, Waterfield J, Bartlam B ve ark.
Saturation in qualitative research: exploring its conceptualization and
operationalization. Quality \& Quantity. 2018;52(4):1893-1907.
\url{https://doi.org/10.1007/s11135-017-0574-8}

\bibitem[\citeproctext]{ref-sharpe2002siblings}
Sharpe D, Rossiter L. Siblings of Children with a Chronic Illness: A
Meta-Analysis. Journal of Pediatric Psychology. 2002;27(8):699-710.
\url{https://doi.org/10.1093/jpepsy/27.8.699}

\bibitem[\citeproctext]{ref-silina2023mediating}
Silina E, Taube M, Zolovs M. Exploring the Mediating Role of Parental
Anxiety in the Link between Children's Mental Health and Glycemic
Control in Type 1 Diabetes. International Journal of Environmental
Research and Public Health. 2023;20(19):6849.
\url{https://doi.org/10.3390/ijerph20196849}

\bibitem[\citeproctext]{ref-silva2015disagreement}
Silva N, Crespo C, Carona C, Bullinger M, Canavarro MC. Why the
(dis)agreement? Family context and child--parent perspectives on
health-related quality of life and psychological problems in paediatric
asthma. Child: Care, Health and Development. 2015;41(1):112-121.
\url{https://doi.org/10.1111/cch.12147}

\bibitem[\citeproctext]{ref-tan2024siblingreview}
Tan RB, Chan PY, Shorey S. Experiences of siblings of children with
chronic pediatric conditions: a qualitative systematic review. European
Journal of Pediatrics. 2024;184(1):44.
\url{https://doi.org/10.1007/s00431-024-05826-7}

\bibitem[\citeproctext]{ref-taylorDeVocht2011separate}
Taylor B, Vocht H de. Interviewing Separately or as Couples?
Considerations of Authenticity of Method. Qualitative Health Research.
2011;21(11):1576-1587. \url{https://doi.org/10.1177/1049732311415288}

\bibitem[\citeproctext]{ref-tong2007coreq}
Tong A, Sainsbury P, Craig J. Consolidated criteria for reporting
qualitative research (COREQ): a 32-item checklist for interviews and
focus groups. International Journal for Quality in Health Care.
2007;19(6):349-357. \url{https://doi.org/10.1093/intqhc/mzm042}

\bibitem[\citeproctext]{ref-tracy2010qualityCriteria}
Tracy SJ. Qualitative {Quality}: {Eight} {``{Big}-{Tent}''} {Criteria}
for {Excellent} {Qualitative} {Research}. Qualitative Inquiry.
2010;16(10):837-851. \url{https://doi.org/10.1177/1077800410383121}

\bibitem[\citeproctext]{ref-williams2009conflict}
Williams LB, Laffel LMB, Hood KK. Diabetes-specific family conflict and
psychological distress in paediatric Type 1 diabetes. Diabetic Medicine.
2009;26(9):908-914.
\url{https://doi.org/10.1111/j.1464-5491.2009.02794.x}

\end{CSLReferences}




\end{document}
